\ExerciseGroupsStart
\ExerciseSection

\begin{enumerate}
\def\labelenumi{\arabic{enumi})}
\tightlist
\item
  A positive real-valued function \(f\) on \(\mathbf{X} \times \mathbf{X}\) is a pseudometric on \(\mathbf{X}\) if and only if \(f(x, x) = 0\) for all \(x \in \mathbf{X}\) and
\end{enumerate}

\[
f (x, y) \leqslant f (x, z) + f (y, z)
\]

for all \(x, y, z\) in \(\mathbf{X}\) .

\begin{enumerate}
\def\labelenumi{\arabic{enumi})}
\setcounter{enumi}{1}
\tightlist
\item
  Let \(f\) be a mapping of \(\mathbf{X} \times \mathbf{X}\) into \([0, +\infty]\) . Then the family of sets \(\overline{f}^{1}([0, a])\) , where \(a\) runs through the set of all real numbers \(>0\) , forms a fundamental system of entourages of a uniformity on \(\mathbf{X}\) if and only if \(f\) satisfies the following conditions: \(a) f(x, x) = 0\) for all \(x \in \mathbf{X}\) ; \(b)\) for each \(a > 0\) there exists \(b > 0\) such that the relation \(f(x, y) \leqslant b\) implies \(f(y, x) \leqslant a\) ; \(c)\) for each \(a > 0\) there exists \(c > 0\) such that the relations \(f(x, z) \leqslant c\) and \(f(z, y) \leqslant c\) imply \(f(x, y) \leqslant a\) .
\end{enumerate}

The conditions \(b)\) and \(c)\) are satisfied in particular if there exists a mapping \(\varphi\) of \(I = [0, +\infty]\) onto itself which is continuous and zero at the point \(0\) , and a mapping \(\psi\) of \(I \times I\) into \(I\) which is continuous and zero at the point \((0, 0)\) , such that we have identically

\[
f (y, x) \leqslant \varphi (f (x, y)) \quad \text { and } \quad f (x, y) \leqslant \psi (f (x, z), f (z, y)).
\]

\begin{enumerate}
\def\labelenumi{\arabic{enumi})}
\setcounter{enumi}{2}
\tightlist
\item
  Let X be a topological space, let C be the topology of X, let \((f_{t})\) be a saturated family (no. 2) of pseudometrics on X, and let U be the uniformity defined by the family \((f_{t})\) .
\end{enumerate}

\begin{enumerate}
\def\labelenumi{\alph{enumi})}
\item
  The topology induced by \(\mathcal{U}\) is coarser than \(\mathfrak{C}\) if and only if the \(f_{t}\) are continuous on \(\mathbf{X} \times \mathbf{X}\) (with respect to the product of the topology \(\mathfrak{C}\) by itself).
\item
  The topology induced by \(\mathcal{U}\) is finer than \(\mathcal{C}\) if and only if, for each \(x_0 \in \mathbf{X}\) and each neighbourhood \(\mathbf{V}\) of \(x_0\) (in the topology \(\mathcal{C}\) ) there is an index \(\iota\) and a real number \(a > 0\) such that \(f_{\iota}(x_0, x) \geqslant a\) for all \(x \in \mathbb{C}\mathbf{V}\) .
\end{enumerate}

\begin{enumerate}
\def\labelenumi{\arabic{enumi})}
\setcounter{enumi}{3}
\tightlist
\item
  Let X be a non-Hausdorff uniformizable space and let R be the relation `` \(y \in \overline{\{x\}}\) '' between two points \(x, y\) of X.
\end{enumerate}

\begin{enumerate}
\def\labelenumi{\alph{enumi})}
\item
  Show that R is an equivalence relation on X, that every continuous mapping of X into a Hausdorff space is compatible with the relation R (Set Theory, R, § 5, no. 7), and that the quotient space X/R is completely regular.
\item
  If \(\mathcal{U}\) is any uniformity compatible with the topology of X, then the Hausdorff uniformity associated with \(\mathcal{U}\) (Chapter II, § 3, no. 9) is defined on the quotient space X/R and is compatible with the topology of X/R. The topological space X/R is called the completely regular space associated with the uniformizable space X.
\end{enumerate}

\begin{enumerate}
\def\labelenumi{\arabic{enumi})}
\setcounter{enumi}{4}
\tightlist
\item
  Let X be a uniformizable space. Show that the family of all pseudometrics on X which are continuous on \(X \times X\) defines a uniformity on X compatible with the topology of X. This uniformity \(U_{0}\) is called the universal uniformity on X; it is the finest of all uniformities compatible with the topology of X. If Y is any uniform space and if \(f: X \to Y\) is a continuous mapping, then f is uniformly continuous with respect to the universal uniformity on X; this is not the case for any other uniformity on X compatible with its topology. If there exists a uniformity U on X, compatible with its topology and such that X, endowed with this uniformity, is a complete space, then X is also a complete space with respect to any uniformity \(U'\) which is finer than U and coarser than \(U_{0}\) . (Observe that the Cauchy filters are the same for \(U'\) as for \(U_{0}\) .)
\end{enumerate}

¶ 6) a) Let X be a completely regular space and let U be a uniformity compatible with the topology of X. Let U* be the coarsest uniformity on X which makes uniformly continuous all the mappings of X into {[}0, 1{]} which are uniformly continuous with respect to U. Show that the uniformity U* is Hausdorff and compatible with the topology of X, and that X is precompact with respect to U*.

\begin{enumerate}
\def\labelenumi{\alph{enumi})}
\setcounter{enumi}{1}
\tightlist
\item
  Let K be a compact space. Show that every mapping \(f: X \to K\) which is uniformly continuous with respect to U is also uniformly continuous with respect to \(U^{*}\) (note that the unique uniformity on K is the coarsest for which all the continuous mappings of K into the interval {[}0, 1{]} are uniformly continuous). Hence show that \(U^{*}\) is the finest of the uniformities on X which are coarser than U and with respect to which X is precompact.
\end{enumerate}

\begin{enumerate}
\def\labelenumi{\arabic{enumi})}
\setcounter{enumi}{6}
\tightlist
\item
  Let X be a completely regular space. If \(U_{0}\) denotes the universal uniformity on X (Exercise 5), then \(U_{0}^{*}\) (Exercise 6) is the coarsest uniformity on X for which all continuous mappings of X into {[}0, 1{]} are uniformly continuous. Let \(\beta X\) denote the compact space obtained by completing X with respect to the uniformity \(U_{0}^{*}\) . \(\beta X\) is called the Stone-Čech compactification of X.
\end{enumerate}

\begin{enumerate}
\def\labelenumi{\alph{enumi})}
\item
  Show that every continuous mapping of X into a compact space K can be extended to a continuous mapping of \(\beta X\) into K (cf.~Set Theory, Chapter IV, § 3).
\item
  Let \(f\) be a homeomorphism of \(\mathbf{X}\) onto a dense subset \(\mathbf{X}'\) of a compact space \(\mathbf{K}\) , and let \(\overline{f}\) be the continuous mapping \(\beta \mathbf{X} \to \mathbf{K}\) which extends \(f\) . Show that \(\overline{f}(\beta \mathbf{X} - \mathbf{X}) = \mathbf{K} - \mathbf{X}'\) . {[}Suppose that there is a point \(a \in \beta \mathbf{X} - \mathbf{X}\) such that \(\overline{f}(a) = f(b)\) , where \(b \in \mathbf{X}\) ; consider a continuous mapping \(g: \beta \mathbf{X} \to [0, 1]\) such that \(g(a) = 1\) and \(g(b) = 0\) , and let \(P\) denote the set of all \(x \in \beta X\) such that \(g(x) \geqslant 1/2\) . There is a continuous mapping \(h: K \to [0, 1]\) such that \(h(f(b)) = 0\) and \(h(f(x)) = 1\) for all \(x \in P\) ; show that the relation \(h(f(a)) = 0\) contradicts the fact that \(a\) is in the closure of \(P\) in \(\beta X\) .{]}
\end{enumerate}

\begin{enumerate}
\def\labelenumi{\arabic{enumi})}
\setcounter{enumi}{7}
\tightlist
\item
  Let X be a completely regular space and let \(\beta X\) be its Stone-Cech compactification.
\end{enumerate}

\begin{enumerate}
\def\labelenumi{\alph{enumi})}
\item
  A filter \(\mathfrak{F}\) on X is said to be completely regular if \(\mathfrak{F}\) has a base \(\mathfrak{B}\) consisting of open sets such that for each set \(A \in \mathfrak{B}\) there exists a set \(B \in \mathfrak{B}\) contained in A, and a continuous mapping \(f: X \to [0, 1]\) which is equal to 0 on B and 1 on \(\mathbb{C}A\) . A completely regular filter is said to be maximal if there exists no strictly finer completely regular filter. Show that, if \(\mathfrak{F}\) is any completely regular filter on E, there is a maximal completely regular filter finer than \(\mathfrak{F}\) (use Zorn's lemma).
\item
  A completely regular filter \(\mathfrak{F}\) is maximal if and only if, for each pair of open sets A, B of X such that \(B \subset A\) and such that there is a continuous mapping of X into \([0, 1]\) which is equal to 0 on B and equal to 1 on \(\complement A\) , we have either \(A \in \mathfrak{F}\) , or else \(A \notin \mathfrak{F}\) there is a set of \(\mathfrak{F}\) which does not meet B (if all the sets of \(\mathfrak{F}\) meet B, consider the filter generated by the sets of \(\mathfrak{F}\) and the sets \(\overline{f}([0, a[)\) , where \(a \in ]0, 1[\) ).
\item
  Show that every maximal completely regular filter \(\mathfrak{F}\) is a Cauchy filter on \(\mathbf{X}\) with respect to the uniformity induced by that of \(\beta\mathbf{X}\) {[}show that each continuous mapping \(f: \mathbf{X} \to [\mathrm{0}, \mathrm{1}]\) has a limit with respect to \(\mathfrak{F}\) , by arguing by reductio ad absurdum and using \(b)\) {]}.
\end{enumerate}

\(d)\) Two distinct maximal completely regular filters \(\mathfrak{F},\mathfrak{F}^{\prime}\) cannot converge to the same point of \(\beta\mathbf{X}\) {[}observe that by \(b)\) there exists an open set \(\mathbf{A}\in\mathfrak{F}\) and an open set \(\mathbf{A}^{\prime}\in\mathfrak{F}^{\prime}\) such that \(\mathbf{A}\cap\mathbf{A}^{\prime}=\varnothing\) ; then argue by contradiction, using axiom \((\mathrm{O}_{\mathbf{IV}})\) {]}.

\begin{enumerate}
\def\labelenumi{\alph{enumi})}
\setcounter{enumi}{4}
\tightlist
\item
  The trace on X of the neighbourhood filter of any point \(x_{0} \in \beta X\) is a maximal completely regular filter B (argue by contradiction, using the definitions of a completely regular filter and of the neighbourhood filter of a point in \(\beta X\)) .
\end{enumerate}

\begin{enumerate}
\def\labelenumi{\arabic{enumi})}
\setcounter{enumi}{8}
\tightlist
\item
  \begin{enumerate}
  \def\labelenumii{\alph{enumii})}
  \tightlist
  \item
    Let X be a topological space, let \(\mathcal{G}\) be its topology and let \((f_{\iota})_{\iota \in \mathbb{I}}\) be a family of mappings of X into K = {[}0, 1{]} which are continuous in the topology \(\mathcal{G}\) . Let \(\mathcal{G}_0\) be the coarsest topology on X for which all the \(f_{\iota}\) are continuous. Show that \(\mathcal{G}_0 = \mathcal{G}\) provided that the family of sets \(\overline{f}_{\iota}^{1}([0, a[)\) , where \(\iota \in I\) and \(a \in ]0, 1[,\) is a subbase of \(\mathcal{G}\) (Chapter I, § 2, no. 3). The space X is then uniformizable, and if it is Hausdorff it is homeomorphic to a subspace of the cube K \(^{\mathbf{I}}\) .
  \end{enumerate}
\end{enumerate}

\begin{enumerate}
\def\labelenumi{\alph{enumi})}
\setcounter{enumi}{1}
\tightlist
\item
  Suppose that the family \((f_{t})\) is the family of all continuous mappings of X into K (with respect to C). Show that, if Y is any compact space, every mapping of X into Y which is continuous with respect to C is continuous with respect to \(C_{0}\) (embed Y in a cube). The topology C is uniformizable if and only if \(C_{0} = C\) .
\end{enumerate}

\begin{enumerate}
\def\labelenumi{\arabic{enumi})}
\setcounter{enumi}{9}
\tightlist
\item
  Let X be a completely regular space, let K be a compact subset of X, and let V be a neighbourhood of K in X.
\end{enumerate}

\begin{enumerate}
\def\labelenumi{\alph{enumi})}
\item
  Show that there is a continuous mapping of X into {[}0, 1{]} which is equal to 1 on K and equal to 0 on \(\complement V\) {[}use \((\mathrm{O}_{\mathbf{IV}})\), covering K by a finite number of suitably chosen neighbourhoods{]}.
\item
  Let \(X'\) be the quotient space of X obtained by identifying all the points of K. Show that \(X'\) is completely regular.
\end{enumerate}

¶ 11) Let X be a completely regular space. Two closed sets A, B in X are said to be completely separated if there exists a continuous mapping f of X into {[}0, 1{]} such that \(f(x) = 0\) on A and \(f(x) = 1\) on B. a) Let \(\beta X\) be the Stone-Čech compactification of X. Show that two closed sets A and B in X are completely separated if and only if their closures in \(\beta X\) do not intersect {[}use Exercises 7 a) and 10 a){]}.

\begin{enumerate}
\def\labelenumi{\alph{enumi})}
\setcounter{enumi}{1}
\tightlist
\item
  Let \(h\) be a homeomorphism of \(\mathbf{X}\) onto a dense subset \(\mathbf{X}'\) of a compact space \(\mathbf{K}\) , and let \(\overline{h}\) be the continuous mapping \(\beta \mathbf{X} \to \mathbf{K}\) which extends \(h\) . Suppose that, for each pair of completely separated closed subsets \(\mathbf{A}, \mathbf{B}\) of \(\mathbf{X}\) , the closures of \(\overline{h}(\mathbf{A})\) and \(h(\mathbf{B})\) in \(\mathbf{K}\) do not intersect. Show that \(\overline{h}\) is a homeomorphism of \(\beta \mathbf{X}\) onto \(\mathbf{K}\) . (Show that \(\overline{h}\) is injective: if \(a, b\) are distinct points of \(\beta \mathbf{X}\) such that \(\overline{h}(a) = \overline{h}(b)\) , consider a continuous mapping \(f: \beta \mathbf{X} \to [\mathbf{0}, \mathbf{1}]\) such that \(f(a) = 0\) , \(f(b) = 1\) , and consider the sets \(X \cap \overline{f}[0, 1/3]\) and \(X \cap \overline{f}[2/3, 1]\) .)
\end{enumerate}

¶ 12) Let X be a completely regular space and let βX be its Stone-Cech compactification.

\begin{enumerate}
\def\labelenumi{\alph{enumi})}
\item
  Let \(f\) be a finite continuous real-valued function on \(\beta X\) , such that \(f(x) > 0\) for all \(x \in X\) , and such that the set \(A = \overline{f}^1(0) \subset \beta X - X\) is not empty. Then there is a sequence \((a_n)\) of points of \(X\) such that the sequence of numbers \(\lambda_n = f(a_n) > 0\) is strictly decreasing and such that \(\lim_{n \to \infty} \lambda_n = 0\) . For each integer \(n > 0\) , let \(I_n\) be an open interval with centre \(\lambda_n\) in \(R_+^*\) , such that the closed intervals \(\overline{I}_n\) are mutually disjoint, and let \(M_n = \overline{f}^1(I_n) \cap X\) . For each subset \(H\) of \(N\) , let \(M_H\) be the union of the \(M_n\) for which \(n \in H\) ; for each filter \(\mathfrak{F}\) on \(N\) which is finer than the Fréchet filter, let \(\mathfrak{F}'\) denote the filter on \(X\) which has the \(M_H(H \in \mathfrak{F})\) as a base. Show that \(\mathfrak{F}'\) is a completely regular filter (Exercise 8), and that if \(\mathfrak{F}_1, \mathfrak{F}_2\) are two distinct ultrafilters on \(N\) , each of which is finer than the Fréchet filter, then there is no filter on \(X\) which is finer than each of \(\mathfrak{F}_1', \mathfrak{F}_2'\) . Deduce that Card \((A) \geqslant 2^{2\text{Card}(N)}\) {[}use Exercise 8 d) and Chapter I, § 4, Exercise 5{]}.
\item
  Suppose that X is infinite and discrete. Show that the uniformity induced on X by that of \(\beta\mathbf{X}\) is the uniformity of finite partitions (Chapter II, § 2, no. 2 and § 4, Exercise 12) {[}use Exercise 11 b){]}. Hence show that Card \((\beta\mathbf{X}) = 2^{2^{\operatorname{Card}(\mathbb{N})}}\) (cf.~Chapter I, § 4, Exercise 5). Show that there is a continuous real-valued function \(f \geqslant 0\) on X such that \(\overline{f}^{1}(0)\) is not empty and is contained in \(\beta\mathbf{X} - \mathbf{X}\) .
\item
  Under the same hypotheses as in b), show that if A is an infinite closed subset of \(\beta X\) , then \(\operatorname{Card}(A) \geqslant 2^{2\operatorname{Card}(N)}\) . {[}Show first that there exists an infinite sequence \((a_{n})\) of points of A, and for each n a neighbourhood \(V_{n}\) of \(a_{n}\) in \(\beta X\) , the \(V_{n}\) being mutually disjoint; deduce that every bounded real-valued function f defined on the set D of the \(a_{n}\) can be extended to a continuous function on \(\beta X\) ; for this, consider a real-valued function g defined on X and equal to \(f(a_{n})\) at each point of \(X \cap V_{n}\) . Hence conclude that \(\overline{D} \subset A\) is homeomorphic to the Stone-Čech compactification of D (cf.~§ 4, Exercise 17).{]}
\end{enumerate}

\begin{enumerate}
\def\labelenumi{\arabic{enumi})}
\setcounter{enumi}{12}
\item
  Let G be a locally compact, non-compact topological group. Show that the uniformity induced on the product space \(H = G \times G\) by the uniformity of its Stone-Čech compactification \(\beta H\) is strictly finer than the uniformity induced by that of \((\beta G) \times (\beta G)\) . {[}Assuming the result false, show, by applying Exercise 7 a) to the mapping \((x, y) \to xy^{-1}\) of H into G, that G would be isomorphic to a subgroup of a compact group; now apply Chapter III, § 3, no. 3, Proposition 4, Corollary 1.{]}
\item
  If a topological space X is such that every point of X has a closed neighbourhood which is a uniformizable subspace of X, show that X is uniformizable.
\end{enumerate}

¶ 15) a) Let X be a locally compact space and let \(\Phi\) be the set of all continuous mappings \(g: \mathbf{X} \to [\mathfrak{0}, \mathfrak{1}]\) which are zero on the complement of some compact set (depending on \(g\) ). Let \(\mathcal{U}_1\) be the coarsest uniformity on X for which all the mappings \(g \in \Phi\) are uniformly continuous. Show that \(\mathcal{U}_1\) is compatible with the topology of X, that the completion \(\hat{\mathbf{X}}\) of X with respect to \(\mathcal{U}_1\) is compact, and that \(\mathbf{X} - \hat{\mathbf{X}}\) consists of at most one point (show that if X is not compact, the filter of complements of relatively compact subsets of X is a Cauchy filter with respect to \(\mathcal{U}_1\) ).

\begin{enumerate}
\def\labelenumi{\alph{enumi})}
\setcounter{enumi}{1}
\item
  Show that \(\mathfrak{U}_1\) is the coarsest uniformity compatible with the topology of \(\mathbf{X}\) .
\item
  Conversely, let X be a completely regular space such that the set of uniformities compatible with the topology of X has a coarsest element \(U_{1}\) . Show that X is locally compact. (Using Exercise 6, show that X is precompact with respect to \(U_{1}\) ; then show that the complement of X in its completion with respect to \(U_{1}\) cannot have more than one point.)
\item
  Suppose that X is locally compact and \(\sigma\) -compact, and therefore the union of an increasing sequence \((\mathbf{U}_n)\) of relatively compact open sets such that \(\overline{\mathbf{U}}_n \subset \mathbf{U}_{n+1}\) (Chapter I, § 9, no. 9, Proposition 15). Show that there exists a continuous real-valued function \(f\) on \(\mathbf{X}\) , such that \(f(x) \leqslant n\) for \(x \in \overline{\mathbf{U}}_n\) and \(f(x) \geqslant n\) for \(x \in \mathbb{C} \overline{\mathbf{U}}_n\) {[}cf.~Exercise 10 a){]}. Let \(\mathcal{U}_2\) be the coarsest uniformity on \(\mathbf{X}\) for which \(f\) and all the functions \(g \in \Phi\) are uniformly continuous. Show that \(\mathcal{U}_2\) is compatible with the topology of \(\mathbf{X}\) and that there is an entourage \(\mathbf{V}\) of \(\mathcal{U}_2\) such that \(\mathbf{V}(x)\) is relatively compact in \(\mathbf{X}\) for all \(x \in \mathbf{X}\) (cf.~Chapter II, § 4, Exercise 9).
\end{enumerate}

¶ 16) Let X be a complete Hausdorff uniform space, and let U be the uniformity of X.

\begin{enumerate}
\def\labelenumi{\alph{enumi})}
\tightlist
\item
  Let A be a subset of X which is the union of a sequence \((\mathbf{F}_{n})\) of closed sets, and let \(x \notin A\) . Show that there is a continuous finite real-valued function \(f \geqslant 0\) on X, such that \(f(x) = 0\) and \(f(y) > 0\) for all \(y \in A\) . (If \(g_{n}\) is a continuous mapping of X into {[}0, 1{]} such that \(g_{n}(x) = 0\) and \(g_{n}(y) = 1\) for all \(y \in F_{n}\) , consider the function
\end{enumerate}

\[
f (x) = \sum_ {n = 0} ^ {\infty} g _ {n} (x) / 2 ^ {n}.
\]

\begin{enumerate}
\def\labelenumi{\alph{enumi})}
\setcounter{enumi}{1}
\tightlist
\item
  Let \(h \geqslant 0\) be a continuous real-valued function on \(\mathbf{X}\) , and let \(\mathbf{V}\) be the open set \(\overline{h}^{1}(]0, +\infty[)\) . Let \(\mathcal{U}'\) be the uniformity on \(\mathbf{V}\) which is the least upper bound of the uniformity induced by U and the uniformity defined by the pseudometric
\end{enumerate}

\[
r (x, y) = \left| \frac {\mathrm{1}}{h (x)} - \frac {\mathrm{1}}{h (y)} \right|.
\]

Show that V is complete with respect to \(U'\) .

\begin{enumerate}
\def\labelenumi{\alph{enumi})}
\setcounter{enumi}{2}
\tightlist
\item
  Let B be a subset of X which is the intersection of a family of sets \((\mathbf{A}_{\lambda})\) , each of which is the union of a countable family of closed sets. Show that there exists a uniformity on B which is compatible with the topology induced on B by the topology of X, and such that B is complete with respect to this uniformity {[}use b){]}.
\end{enumerate}

\begin{enumerate}
\def\labelenumi{\arabic{enumi})}
\setcounter{enumi}{16}
\item
  Let X be a topological space such that each \(x \in X\) has a fundamental system of neighbourhoods which are both open and closed. Show that X is uniformizable.
\item
  Let X be a countable completely regular space. Show that, for each \(x \in X\) , the neighbourhoods of x which are both open and closed form a fundamental system of neighbourhoods of x (cf.~§ 6, no. 4).
\item
  Let X be a locally compact space. Show that every function \(f \geqslant 0\) which is lower semi-continuous on X is the upper envelope of a family of continuous functions \(\geqslant 0\) , each of which is zero on the complement of a compact set.
\item
  \begin{enumerate}
  \def\labelenumii{\alph{enumii})}
  \tightlist
  \item
    Let X be a completely regular space, and let a be a point of X. Then there exists a continuous function \(f \geqslant 0\) on X such that \(f(a) = 0\) and \(f(x) > 0\) whenever \(x \neq a\) , if and only if there is a sequence \((\mathbf{V}_n)\) of neighbourhoods of a in X such that \(\bigcap_{n} V_n = \{a\}\) {[}cf.~Exercise 16 a){]}.
  \end{enumerate}
\end{enumerate}

\begin{enumerate}
\def\labelenumi{\alph{enumi})}
\setcounter{enumi}{1}
\tightlist
\item
  Let X be a completely regular space, each point of which has a countable fundamental system of neighbourhoods. Show that, in the Stone-Čech compactification \(\beta\mathrm{X}\) , the set X is equal to the set of points which have a countable fundamental system of neighbourhoods in \(\beta\mathrm{X}\) {[}use a) and Exercise 12{]}. Let \(\mathbf{X}'\) be another completely regular space, each point of which has a countable fundamental system of neighbourhoods. Show that if the Stone-Čech compactifications \(\beta\mathrm{X}\) and \(\beta\mathrm{X}'\) are homeomorphic, then X and \(\mathbf{X}'\) are homeomorphic.
\end{enumerate}

¶ 21) a) Let X be a topological space. Show that the following two properties are equivalent: α) every finite continuous real-valued function on X is bounded; β) every bounded continuous real-valued function on X attains its bounds. X is said to be pseudo-compact if it has these properties. If X is pseudo-compact and if f is any continuous mapping of X into a topological space \(X'\) , then \(f(X)\) is pseudo-compact.

\begin{enumerate}
\def\labelenumi{\alph{enumi})}
\setcounter{enumi}{1}
\item
  Consider the following two properties of a topological space X: \(\gamma\) if \((\mathbf{U}_n)\) is any countable open covering of X, then X is the union of a finite number of the closures \(\overline{\mathbf{U}}_n\) ; \(\delta\) every countable filter base on X which is formed of open sets has at least one cluster point. Show that \(\gamma\) and \(\delta\) are equivalent and that they imply that X is pseudo-compact. In particular, every absolutely closed topological space (Chapter I, § 9, Exercise 19) is pseudo-compact. If X has properties \(\gamma\) and \(\delta\) , then so does every subspace of X which is the closure of an open set of X {[}cf.~§ 4, Exercise 26 b){]}.
\item
  Consider the following property of a topological space \(\mathbf{X}\) : \(\zeta)\) every locally finite open covering of \(\mathbf{X}\) is finite. Show that \((\zeta) \Longrightarrow \gamma)\) . {[}If \((\mathbf{U}_n)\) is an increasing sequence of open subsets of \(\mathbf{X}\) which cover \(\mathbf{X}\) and are such that \(\mathbf{U}_{n+1} \notin \overline{\mathbf{U}}_n\) , consider the covering formed by the sets \(\mathbf{U}_{n+1} \cap \mathbb{C}\overline{\mathbf{U}}_n\) and the complement of a sequence \((a_n)\) such that
\end{enumerate}

\[
a _ {n + 1} \in \mathrm{U} _ {n + 1} \cap \left[ \overline {{\mathrm{U}}} _ {n}. \right]
\]

If X is regular, show that \(\gamma\Longrightarrow\zeta\) . {[}Let \((\mathrm{U}_{\alpha})\) be an infinite, locally finite, open covering of X. Define by induction a sequence \((\alpha_{n})\) of indices, a sequence \((x_{n})\) of points of X, and for each \(x_{n}\) two open neighbourhoods \(V_{n}\) and \(W_{n}\) of \(x_{n}\) such that: (i) we have \(\overline{V}_{n}\subset W_{n}\) , \(W_{n}\subset U_{\alpha_{n}}\) , and \(W_{n}\) meets only a finite number of sets \(U_{\alpha}\) ; (ii) \(U_{\alpha_{n}}\) meets none of the \(W_{k}\) with indices k\textless n.~Now consider the covering formed by the \(W_{n}\) and the complement of the union of the \(\overline{V}_{n}\) .{]}

\begin{enumerate}
\def\labelenumi{\alph{enumi})}
\setcounter{enumi}{3}
\item
  In a completely regular space X the properties \(\alpha\) , \(\beta\) , \(\gamma\) , \(\delta\) ) and \(\zeta\) are all equivalent, and are equivalent to the following property: \(\theta\) X is precompact in any uniformity compatible with its topology. {[}To show that \(\alpha) \Longrightarrow \zeta\) ), note that if \((\mathbf{U}_n)\) is a countably infinite, locally finite, open covering of X and if \(a_n \in \mathbf{U}_n\) , then there is a continuous real-valued function \(f\) on X such that \(f(a_n) = n\) for each \(n\) {[}use \((\mathrm{O}_{\mathrm{IV}})]\) . To show that \(\theta) \Longrightarrow \alpha\) ), use Exercise 5 and observe that in a precompact space every uniformly continuous real-valued function is bounded. To show that \(\gamma) \Longrightarrow \theta\) ), note that if X is not precompact with respect to a uniformity \(\mathcal{U}\) , then there is a symmetric entourage V of \(\mathcal{U}\) and a sequence \((x_n)\) of points of X such that no two of the sets \(\mathrm{V}(x_n)\) intersect.{]}
\item
  If X is a completely regular pseudo-compact space which is complete with respect to some uniformity compatible with its topology, then X is compact.
\item
  On a completely regular pseudo-compact space X, the universal uniformity (Exercise 5) is induced by the uniformity of the Stone-Čech compactification of X, and is the unique uniformity, compatible with the topology of X, for which all the continuous mappings of X into {[}0, 1{]} are uniformly continuous.
\end{enumerate}

\begin{enumerate}
\def\labelenumi{\arabic{enumi})}
\setcounter{enumi}{21}
\tightlist
\item
  Let X be a Hausdorff uniform space, let Y be a closed subset of X, and let f be a bounded, uniformly continuous, real-valued function on Y. Show that f has an extension \(\overline{f}\) to X which is bounded and uniformly continuous on X. {[}We may assume that \(f(Y) \subset [0, 1]\) . For each dyadic number \(r = k/2^{n} \in [0, 1]\) , let \(A(r)\) be the set of all \(x \in Y\) such that \(f(x) \leqslant r\) , and let \(B(r) = A(r) \cup (X - Y)\) . Define by induction * (as in the proof of Theorem 1 of § 4, no. 1) * an open set \(U(r)\) for each dyadic number \(r \in [0, 1]\) such that, whenever r and \(r'\) are two dyadic numbers and \(r < r'\) , there is an entourage V of the uniformity for X for which
\end{enumerate}

\[
\mathbf {V} (\mathbf {A} (r)) \subset \mathbf {U} (r ^ {\prime}), \quad \mathbf {V} (\mathbf {U} (r)) \subset \mathbf {U} (r ^ {\prime}), \quad \mathbf {V} (\mathbf {U} (r)) \subset \mathbf {B} (r ^ {\prime}).
\]

Then define \(\overline{f}(x)\) to be the greatest lower bound of the dyadic numbers \(r\) such that \(x \in \mathrm{U}(r).]\)

\ExerciseSection

\begin{enumerate}
\def\labelenumi{\arabic{enumi})}
\item
  Let X be a Hausdorff uniform space and let \((f_{t})\) be a family of pseudometrics on X which define the uniformity of X. Let \(X_{t}\) be the metric space associated with the uniform space obtained by endowing X with the structure defined by the single pseudometric \(f_{t}\) (no. 1). Show that X is isomorphic to a subspace of the product uniform space \(\prod_{i} X_{i}\) .
\item
  In a connected metric space X for which the metric is not bounded on \(X \times X\) , show that a sphere cannot be empty.
\end{enumerate}

¶ 3) a) Let X be a compact metric space. If the closure in X of every open ball is the closed ball of the same centre and radius, then every ball in X is a connected set {[}if x and y are two points of X, and if S is the closed ball with centre x and radius \(d(x, y)\) , show that for each \(\varepsilon > 0\) the set \(A_{y,\varepsilon}\) of points of S which can be joined to y by a \(V_{\varepsilon}\) -chain contained in S (Chapter II, § 4, no. 4) contains points z such that \(d(x, z) < d(x, y)\) {]}. Deduce that \(x \in A_{y,\varepsilon}\) for all \(\varepsilon > 0\) , arguing by contradiction and using Weierstrass's theorem (Chapter IV, § 6, no. 1, Theorem 1).

\begin{enumerate}
\def\labelenumi{\alph{enumi})}
\setcounter{enumi}{1}
\tightlist
\item
  In the space \(\mathbf{Y} = \mathbf{R}^2\) with the metric
\end{enumerate}

\[
d (x, y) = \max \left(\left| x _ {1} - y _ {1} \right|, \left| x _ {2} - y _ {2} \right|\right),
\]

let X be the compact subspace consisting of all points \((x_{1}, x_{2})\) such that either \(x_{1} = 0\) and \(0 \leqslant x_{2} \leqslant 1\) , or \(0 \leqslant x_{1} \leqslant 1\) and \(x_{2} = 0\) . Show that every ball in X is connected but that the closure of an open ball is not necessarily the closed ball with the same centre and radius.

\begin{enumerate}
\def\labelenumi{\arabic{enumi})}
\setcounter{enumi}{3}
\tightlist
\item
  A metric space \(\mathbf{X}\) is said to be an ultrametric space if its metric \(d\) satisfies the inequality
\end{enumerate}

\[
d (x, y) \leqslant \sup \left(d (x, z), d (y, z)\right)
\]

for all \(x, y, z\) in \(\mathbf{X}\) (which implies the triangle inequality) (cf.~§ 6, Exercise 2).

\begin{enumerate}
\def\labelenumi{\alph{enumi})}
\item
  If \(d(x, z) \neq d(y, z)\) , show that \(d(x, y) = \sup (d(x, z), d(y, z))\) .
\item
  Let \(\mathrm{V}_{r}(x)\) be the open ball with centre x and radius r. Show that \(\mathrm{V}_{r}(x)\) is both open and closed in X (and hence that X is totally disconnected) and that for each \(y \in \mathrm{V}_{r}(x)\) we have \(\mathrm{V}_{r}(y) = \mathrm{V}_{r}(x)\) .
\item
  Show that the closed ball \(\mathrm{W}_{r}(x)\) with centre x and radius r is both open and closed in X, and that for each \(y \in \mathrm{W}_{r}(x)\) we have
\end{enumerate}

\[
\mathrm{W} _ {r} (y) = \mathrm{W} _ {r} (x).
\]

The distinct open balls of radius r contained in \(W_{r}(x)\) form a partition of \(W_{r}(x)\) , and the distance between any pair of these balls is equal to r. d) If two balls (open or closed) of X intersect, then one is contained in the other.

\begin{enumerate}
\def\labelenumi{\alph{enumi})}
\setcounter{enumi}{4}
\item
  A sequence \((x_{n})\) of points of X is a Cauchy sequence if and only if \(d(x_{n}, x_{n+1})\) tends to 0 as \(n\) tends to infinity.
\item
  If X is compact, show that for each \(x_{0} \in X\) the set of values of \(d(x_{0}, x)\) in X is a (finite or infinite) countable subset of \([0, +\infty]\) , of which all the points, with the possible exception of 0, are isolated {[}for each value r taken by \(d(x_{0}, x)\) , consider the least upper bound of \(d(x_{0}, x)\) on the set of points where \(d(x_{0}, x) < r\) , and its greatest lower bound on the set of points where \(d(x_{0}, x) > r\) {]}.
\end{enumerate}

\begin{enumerate}
\def\labelenumi{\arabic{enumi})}
\setcounter{enumi}{4}
\tightlist
\item
  Let X be a metric space, d its metric. For each pair \((x, y)\) of points of X, let \(d_{0}(x, y)\) denote the greatest lower bound of the real numbers \(\alpha > 0\) such that x and y can be joined by a \(V_{\alpha}\) -chain (Chapter II, § 4, no. 4). Show that \(d_{0}\) is a pseudometric on X, and that the metric space associated with the uniform space defined by the pseudometric \(d_{0}\) on X (no. 1) is an ultrametric space (Exercise 4).
\end{enumerate}

§ 6) Let X be a metric space. If A and B are two non-empty subsets of X, put

\[
\rho (\mathrm{A}, \mathrm{B}) = \sup _ {x \in \mathbf {A}} d (x, \mathrm{B}) \quad \text { and } \quad \sigma (\mathrm{A}, \mathrm{B}) = \sup \left(\rho (\mathrm{A}, \mathrm{B}), \rho (\mathrm{B}, \mathrm{A})\right);
\]

also put \(\sigma(\emptyset, \emptyset) = 0\) , \(\sigma(\emptyset, A) = \sigma(A, \emptyset) = +\infty\) for any non-empty subset \(A\) of \(X\) . Show that \(\sigma\) is a pseudometric on the set \(\mathfrak{P}(X)\) of all subsets of \(X\) , and that the uniformity defined by \(\sigma\) coincides with the uniformity constructed from that of \(X\) by the procedure of Chapter II, § 1, Exercise 5.

Suppose that X is bounded. Then \(\sigma\) is a metric on the set \(\mathfrak{F}(\mathbf{X})\) of non-empty closed subsets of X. If, moreover, X is complete, show that \(\mathfrak{F}(\mathbf{X})\) is a complete metric space {[}let \(\Phi\) be a Cauchy filter on \(\mathfrak{F}(\mathbf{X})\) ; for each set \(\mathcal{X} \in \Phi\) , let \(S(\mathcal{X})\) be the union of the subsets A of X which belong to \(\mathcal{X}\) ; show that the sets \(S(\mathcal{X})\) form a filter base on X as \(\mathcal{X}\) runs through \(\Phi\) , that this filter base has a non-empty set C of cluster points, and that \(\Phi\) converges to C{]}.

\begin{enumerate}
\def\labelenumi{\arabic{enumi})}
\setcounter{enumi}{6}
\item
  Let X be an infinite discrete space. Show that the uniformity of finite partitions (Chapter II, § 2, no. 2) on X, which is compatible with the topology of X, is not metrizable (otherwise there exists a sequence \((\mathfrak{F}_n)\) of finite partitions of X such that every finite partition of X is formed of sets, each of which is a union of sets belonging to one of the partitions \(\mathfrak{F}_n\) ; from this it would follow that the set of finite partitions of X is countable).
\item
  Let \((\mathbf{X}_{t})\) be an uncountable family of Hausdorff topological spaces, each of which has at least two distinct points. Show that, in the product space \(\prod_{i} X_{i}\) , no point has a countable fundamental system of neighbourhoods.
\item
  Let X be a topological space, each point of which has a countable fundamental system of neighbourhoods.
\end{enumerate}

\begin{enumerate}
\def\labelenumi{\alph{enumi})}
\item
  X is Hausdorff provided that every convergent sequence in X has only one limit.
\item
  If a is a cluster point of an infinite sequence \((x_{n})\) of points of X, then there is an infinite subsequence of \((x_{n})\) which converges to a.
\item
  Let A be a non-empty subset of X, let \(x_{0}\) be a point of \(\overline{A}\) , and let f be a mapping of A into a Hausdorff topological space \(X'\) . Then a point \(a \in X'\) is a limit of f at the point \(x_{0}\) , relative to A, provided that, for every sequence \((x_{n})\) of points of A which converges to \(x_{0}\) , the sequence \((f(x_{n}))\) converges to a.
\item
  With the notation of c), suppose also that every point of \(X'\) has a countable fundamental system of neighbourhoods. If \(a \in X'\) is a cluster point of f at \(x_{0}\) relative to A, then there is a sequence \((x_{n})\) of points of A which converges to \(x_{0}\) and is such that the sequence \((f(x_{n}))\) converges to a.
\end{enumerate}

¶ 10) Let X be a compact space. If there exists a continuous real-valued function f on X×X such that the relation \(f(x,y)=0\) is equivalent to x=y, then X is metrizable. (Show that if V runs through a fundamental system of neighbourhoods of 0 in R, the sets \(\overline{f}^{1}(V)\) form a fundamental system of neighbourhoods of the diagonal \(\Delta\) in X×X.)

¶ 11) Let X be a compact metric space, with metric d.~Show that if f is a mapping of X into X such that, for each pair x, y of points of X, we have \(d(f(x), f(y)) \geqslant d(x, y)\) , then f is an isometry of X onto itself. {[}Let a, b be any two points of X, and put \(f^n = f^{n-1} \circ f\) , \(a_n = f^n(a)\) , \(b_n = f^n(b)\) ; show that, for each \(\varepsilon > 0\) , there exists an index k such that \(d(a, a_k) \leqslant \varepsilon\) and \(d(b, b_k) \leqslant \varepsilon\) , by choosing suitable subsequences of \((a_n)\) and \((b_n)\) ; hence show that \(d(a_1, b_1) = d(a, b)\) and that \(f(\mathbf{X})\) is dense in X.{]}

¶ 12) Let X be a compact metrizable space.

\begin{enumerate}
\def\labelenumi{\alph{enumi})}
\item
  Show that there is a continuous mapping of Cantor's triadic set K (Chapter IV, § 2, no. 5) onto X (cf.~Chapter IV, § 8, Exercise 11).
\item
  If in addition X is totally disconnected and has no isolated points, then X is homeomorphic to K. {[}Argue as in Chapter IV, § 8, Exercise 12, using the fact that every neighbourhood of a point \(x \in \mathbf{X}\) contains a neighbourhood of \(x\) which is both open and closed (Chapter II, § 4, no. 4, Corollary to Proposition 6){]}.
\end{enumerate}

¶ 13) a) On the set R of real numbers, consider the topology C defined as follows: for each y \textgreater{} 0, \(\mathrm{U}_{y}(x)\) denotes the union of the intervals \([x, x + y[ \text{ and } ] - x - y, -x[; C \text{ is the topology for which the } \mathrm{U}_{y}(x) \text{ form a fundamental system of neighbourhoods of } x \text{ as } y \text{ runs through the set of real numbers } > 0. \text{ Let X be the space obtained by endowing the interval } [-1, +1] \text{ with the topology induced by C. Show that X is compact. (Consider a filter F on X; if } x \in X \text{ is a cluster point of F with respect to the topology of the real line, show that either } x \text{ or } -x \text{ is a cluster point of F with respect to C.)}\)

\begin{enumerate}
\def\labelenumi{\alph{enumi})}
\setcounter{enumi}{1}
\item
  Every point of X has a countable fundamental system of neighbourhoods, and X has a countable dense subset, but X is not metrizable (show that its topology has no countable base).
\item
  Let A be an open subset of X. Show that A is the union of a countable family \((I_{n})\) of open intervals contained in {[}0, 1{]}, the intervals \(-I_{n}\) , a subset of the set of left-hand points of the intervals \(I_{n}\) and \(-I_{n}\) , and possibly the point +1. Hence show that A is the union of a countable family of closed subsets of X.
\item
  Let Y be the set \(I \times \{1, 2\}\) , where I is the interval {[}0, 1{]} of R; let \(X_{i}\) denote \(I \times \{i\} (i = 1, 2)\) and let f denote the bijection \((x, 1) \to (x, 2)\) of \(X_{1}\) onto \(X_{2}\) . For each \(x \in I\) , let \(\mathfrak{B}((x, 2))\) denote the set of subsets consisting only of \(\{(x, 2)\}\) , and let \(\mathfrak{B}((x, 1))\) denote the set of subsets of the form \(V_{1} \cup (f(V_{1}) - \{x, 2\})\) , where \(V_{1} = V \times \{1\}\) and V runs through a fundamental system of neighbourhoods of x in I. Show that for each \(y \in Y\) , \(\mathfrak{B}(y)\) is a fundamental system of neighbourhoods of y for a topology on Y. Endowed with this topology, Y is compact and every point of Y has a countable fundamental system of neighbourhoods, but Y has no countable dense subset. Every closed subset of Y contained in \(X_{2}\) is finite; hence show that the compact subset \(X_{1}\) of Y has no countable fundamental system of neighbourhoods. If R is the equivalence relation on Y whose classes are the set \(X_{1}\) and the points of \(Y - X_{1}\) , R is closed and every equivalence class mod R is compact, but there is a point in Y/R which has no countable fundamental system of neighbourhoods {[}cf.~§4, Exercise 24a){]}.
\end{enumerate}

\begin{enumerate}
\def\labelenumi{\arabic{enumi})}
\setcounter{enumi}{13}
\tightlist
\item
  \begin{enumerate}
  \def\labelenumii{\alph{enumii})}
  \tightlist
  \item
    Let X be an accessible topological space (Chapter I, § 8, Exercise 1). Show that the following properties are equivalent:
  \end{enumerate}
\end{enumerate}

α) Every sequence of points of X has a cluster point.

\(\beta)\) No infinite discrete subspace of \(\mathbf{X}\) is closed.

\(\gamma)\) Every countable open covering of \(\mathbf{X}\) contains a finite open covering \(\mathbf{X}\) .

δ) Given any infinite open covering R of X, there exists an open covering S ⊂ R of X, distinct from R.

A topological space X is said to be countably compact if it is Hausdorff and has the above properties.

\begin{enumerate}
\def\labelenumi{\alph{enumi})}
\setcounter{enumi}{1}
\item
  A sequence of points in a countably compact space is convergent if and only if it has a unique cluster point.
\item
  Every closed subspace of a countably compact space is countably compact. Conversely, if X is Hausdorff and if every point of X has a countable fundamental system of neighbourhoods, then every countably compact subspace of X is closed in X.
\item
  Let \(f\) be a continuous mapping of a countably compact space \(\mathbf{X}\) into a Hausdorff space \(\mathbf{X}'\) . Then \(f(\mathbf{X})\) is a countably compact subspace of \(\mathbf{X}'\) .
\item
  Let X be a countably compact space, every point of which has a countable fundamental system of neighbourhoods. Then every sequence of points of X has a convergent subsequence.
\end{enumerate}

\(f\) ) Let \((\mathbf{X}_n)\) be a countable sequence of topological spaces, in each of which every point has a countable fundamental system of neighbourhoods.

The product space \(\prod_{n=1}^{\infty} \mathbf{X}_n\) is then countably compact if and only if each of the spaces \(\mathbf{X}_n\) is countably compact. {[}Use e) and Chapter I, §6, Exercise 16.{]}

\begin{enumerate}
\def\labelenumi{\alph{enumi})}
\setcounter{enumi}{6}
\item
  A countably compact space in which every point has a countable fundamental system of neighbourhoods is regular.
\item
  If a countably compact space has a countable base of open sets, it is compact and therefore metrizable.
\item
  Every lower semi-continuous real-valued function on a countably compact space attains its greatest lower bound. In particular, every countably compact space is pseudo-compact (§ 1, Exercise 21) {[}cf.~§ 4, Exercise 26 b){]}.
\item
  Show that the property \(\delta\) of \(a\) implies the following:
\end{enumerate}

ζ) Every point-finite (§ 4, no. 3) open covering of X contains a finite open covering of X. (Argue by contradiction.) Conversely, a regular space which satisfies ζ) is countably compact {[}show that ζ) then implies β){]}.

¶ 15) Let \(\mathbf{X} = [a, b[\) be the locally compact space defined in Chapter I, § 9, Exercise 12.

\begin{enumerate}
\def\labelenumi{\alph{enumi})}
\item
  A subset of X is relatively compact if and only if it is bounded. Hence show that X is countably compact and therefore (Proposition 15) non-metrizable (observe that every countable subset of X is bounded).
\item
  Show that every point of X has a metrizable neighbourhood (use Proposition 16).
\item
  If A and B are two non-compact closed sets in X, their intersection is not empty (form an increasing sequence in which the even-numbered terms are points of A and the odd-numbered terms are points of B). Deduce that every neighbourhood of a non-compact closed set is the complement of a relatively compact set. If A is the set of non-isolated points of X, show that A is closed and is not the intersection of any countable family of open sets.
\item
  Let \(\mathbf{X}'\) denote the interval \([a, b]\) endowed with the following topology: for each \(x \in \mathbf{X}\) the intervals \(]y, x]\) , where \(y\) runs through the set of elements \(< x\) , form a fundamental system of neighbourhoods of \(x\) ; a fundamental system of neighbourhoods of \(b\) is formed by the sets \(\mathbf{V}_x\) where, for each \(x \in \mathbf{X}\) , \(\mathbf{V}_x\) denotes the union of \(\{b\}\) and the set of points of \([x, b]\) which have a predecessor (i.e., which are isolated in \(\mathbf{X}\) ). Show that \(\mathbf{X}'\) is countably compact but not regular {[}cf.~Exercise 14 g){]} and that the subspace \(\mathbf{X}\) of \(\mathbf{X}'\) is countably compact but not closed {[}cf.~Exercise 14 c){]}.
\end{enumerate}

\begin{enumerate}
\def\labelenumi{\arabic{enumi})}
\setcounter{enumi}{15}
\tightlist
\item
  Let X and Y be two countably compact spaces. Show that the product \(X \times Y\) is countably compact in each of the following two cases: (i) one of X, Y is compact; (ii) one of X, Y is such that each point has a countable fundamental system of neighbourhoods. {[}In case (i), supposing that Y is compact, consider an increasing sequence \((\mathbf{G}_{n})\) of open sets whose union is \(X \times Y\) ; for each n, let \(H_{n}\) be the set of all \(x \in X\) such that \(\{x\} \times Y \subset G_{n}\) ; show that \(H_{n}\) is open and that \(X = \bigcup_{n} H_{n} \cdot\) {]}
\end{enumerate}

¶ 17) Let βN be the Stone-Cech compactification of the discrete space N (§ 1, Exercise 7).

\begin{enumerate}
\def\labelenumi{\alph{enumi})}
\item
  Show that there are two countably compact subspaces A, B of \(\beta N\) such that \(A \cap B = N\) and \(A \cup B = \beta N\) . {[}Let \(\aleph_{\alpha} = 2^{2^{\operatorname{Card}(N)}}\) ; by virtue of § 1, Exercise 12 b), there is a bijection \(\xi \to S_{\xi}\) of the ordinal \(\omega_{\alpha}\) (Set Theory, Chapter III, § 6, Exercise 10) onto the set of countably infinite subsets of X. By transfinite induction, define two injective mappings \(\xi \to x_{\xi}, \xi \to y_{\xi}\) of \(\omega_{\alpha}\) into \(\beta N - N\) such that \(x_{\xi} \in \overline{S}_{\xi}, y_{\xi} \in \overline{S}_{\xi}\) for each \(\xi\) , and such that if P (resp. Q) is the set of the \(x_{\xi}\) (resp. \(y_{\xi}\) ) we have \(P \cap Q = \varnothing\) , \(P \cup Q = \beta N - N\) ; for this, make use of § 1, Exercise 12 b) and c). Then show that A = P ∪ N and B = Q ∪ N satisfy the conditions of the question.{]}
\item
  Show that the product space \(A \times B\) is not countably compact (note that the intersection of \(A \times B\) with the diagonal of \(\beta N \times \beta N\) is an infinite discrete closed subspace of \(A \times B\) ).
\end{enumerate}

¶ 18) Let X be a metric space such that, for each \(x \in X\) , there is an open ball with centre \(x\) which, when considered as a subspace of X, has a countable base. Let \(r_x\) be the least upper bound of the radii of balls with centre \(x\) which have this property.

\begin{enumerate}
\def\labelenumi{\alph{enumi})}
\item
  Show, using Zorn's lemma, that there is a maximal family \((\mathbf{B}_{\alpha})\) of open balls in X which are pairwise disjoint and such that, if \(x_{\alpha}\) is the centre of \(B_{\alpha}\) , the radius of \(B_{\alpha}\) is \(<r_{x_{\alpha}}\) . Show that the union of the \(B_{\alpha}\) is dense in X, and hence that there is a dense subset M of X such that every open ball with centre at an arbitrary point \(x \in X\) and radius \(<r_{x}\) contains only a countable infinity of points of M (note that such a ball can meet only a countable infinity of pairwise disjoint open sets, and use Proposition 12).
\item
  For each point \(x \in \mathbf{M}\) , let \(\mathbf{S}_x\) denote the open ball with centre \(x\) and radius \(\frac{1}{3} r_x\) . Show that the \(\mathbf{S}_x\) cover \(\mathbf{X}\) , and that an arbitrary point of \(\mathbf{X}\) belongs only to a countable infinity of the \(\mathbf{S}_x\) {[}note that if \(y \in \mathbf{S}_x\) , we have \(d(x, y) \leqslant \frac{1}{2} r_y\) {]}.
\item
  Let R be the following equivalence relation between points \(x, y\) of X: there exists a sequence \((z_i)_{1 \leqslant i \leqslant n}\) of points of M such that \(x \in S_{z_i}\) , \(y \in S_{z_n}\) and \(S_{z_i}\) meets \(S_{z_{i+1}}\) for \(1 \leqslant i \leqslant n - 1\) . Show that the equivalence classes (mod R) are subspaces of X which are both open and closed in X and which have a countable base; in other words, X is the topological sum of metric spaces with countable bases (Chapter I, § 2, no. 4).
\end{enumerate}

¶ 19) In a metric space every relatively compact set is bounded. If X is a topological space, show that the following conditions on X are equivalent: α) there exists a metric on X, compatible with its topology, such that every bounded subset of X (with respect to this metric) is relatively compact; β) X is locally compact and has a countable base. {[}To show that α) → β), show that if every bounded set is relatively compact, then X is locally compact and σ-compact. To show that β) → α), note that X is metrizable by the corollary to Proposition 12; if d is a metric on X compatible with the topology of X, and if f is the function defined in § 1, Exercise 15 d), take the uniformity on X defined by the two pseudometrics d(x, y) and \textbar f(x) --- f(y)\textbar.{]}

¶ 20) a) Give an example of a closed equivalence relation R on a metrizable locally compact space X which has a countable base such that X/R is paracompact but such that there exists a point of X/R which has no countable fundamental system of neighbourhoods {[}cf.~Chapter I, § 10, Exercise 17 and Chapter IX, § 4, Exercise 24 a){]}.

\begin{enumerate}
\def\labelenumi{\alph{enumi})}
\setcounter{enumi}{1}
\tightlist
\item
  Let X be a metrizable space and let R be a closed equivalence relation on X. Show that if every point of X/R has a countable fundamental system of neighbourhoods, then each equivalence class (mod R) has a compact frontier in X. (Show that if the result were false there would exist a sequence in X, all of whose points were distinct, which had no cluster point, and whose image in X/R was a sequence which converged to a point distinct from all the points of the sequence.) For each \(z \in X/R\) , let \(C_{z}\) be the inverse image of z in X, and let \(F_{z}\) be the frontier of \(C_{z}\) if this frontier is not empty; if \(C_{z}\) is both open and closed, let \(F_{z}\) be any subset of \(C_{z}\) consisting of a single point. Let Y be the union of the sets \(F_{z}\) as z runs through X/R. Show that \(Y/R_{y}\) is homeomorphic to X/R.
\end{enumerate}

¶ 21) a) Let X be a metrizable space, let d be a bounded metric compatible with the topology of X, let σ be the metric corresponding to d on the set ℱ(X) of non-empty closed subsets of X (Exercise 6); and let R be an open and closed equivalence relation on X. Show that the restriction of σ to X/R is a metric compatible with the quotient topology. {[}Use Exercise 20 b) to show that every class mod R is open and compact. Then prove that if a sequence (zn) in X/R tends to a point a and if φ: X → X/R is the canonical mapping, we have σ(φ̄¹(a), φ̄¹(zₙ)) → 0; argue by contradiction, using the fact that R is open.{]}

\begin{enumerate}
\def\labelenumi{\alph{enumi})}
\setcounter{enumi}{1}
\tightlist
\item
  On the compact interval I = {[}0, 1{]} of R, let R be the equivalence relation whose classes are the points of the Cantor set K (Chapter VI, § 2, no. 5), other than the end-points of the intervals contiguous to K, and the closures of the intervals contiguous to K. Show that the quotient space I/R is homeomorphic to I {[}Chapter IV, § 8, Exercise 16 b){]} but that the metric \(\sigma\) is not compatible with the topology of I/R.
\end{enumerate}

\begin{enumerate}
\def\labelenumi{\arabic{enumi})}
\setcounter{enumi}{21}
\tightlist
\item
  Let X be a Hausdorff topological space with a countable base \((\mathbf{U}_{n})\) , and let R be a Hausdorff equivalence relation on X such that every point of X/R has a countable fundamental system of neighbourhoods. Show that the topology of X/R has a countable base. {[}Let \(\varphi: X \to X/R\) be the canonical mapping, and show that the interiors of finite unions of sets \(\varphi(\mathbf{U}_{n})\) form a base of the topology of X/R: if V is a neighbourhood of a point \(z \in X/R\) , and \((\mathbf{W}_{k})\) is a sequence of sets belonging to the base \((\mathbf{U}_{n})\) , contained in \(\overline{\varphi}^{1}(\mathbf{V})\) and covering \(\overline{\varphi}^{1}(z)\) , show that there is a finite number of indices k such that the union of the \(\varphi(\mathbf{W}_{k})\) is a neighbourhood of z. Argue by contradiction: if this statement were false we could construct a sequence \((y_{n})\) of distinct points of X/R, tending to z and such that \(y_{n}\) did not belong to the union of the \(\varphi(\mathbf{W}_{k})\) for \(k \leqslant n\) ; then show that the union of the \(\overline{\varphi}^{1}(y_{n})\) would be closed in X.{]}
\end{enumerate}

Show that the same conclusion is valid if R is assumed to be closed and every equivalence class mod R is compact (similar method).

¶ 23) A topological space is said to be submetrizable if its topology is finer than the topology of a metrizable space. A submetrizable space is Hausdorff but not necessarily regular (Chapter I, § 8, Exercise 20).

\begin{enumerate}
\def\labelenumi{\alph{enumi})}
\item
  A completely regular space X is submetrizable if and only if there exists a uniformity on X, compatible with the topology of X and defined by a family \(\Phi\) of pseudometrics which contains at least one metric d (cf.~§1, Exercise 3). Let \(\hat{X}\) be the completion of X with respect to this uniformity; the metric d extends to a pseudometric \(\overline{d}\) on \(\hat{X}\) ; show that for each \(x \in \hat{X}\) there is at most one point \(y \in X\) such that \(\overline{d}(x, y) = 0\) .
\item
  Show that if X is a completely regular submetrizable space, there exists a uniformity, compatible with the topology of X, and with respect to which X is complete {[}use a) and § 1, Exercise 16 c){]}. Hence show that if in addition X is pseudo-compact (§ 1, Exercise 21), then X is compact.
\item
  Show that a locally compact, paracompact, submetrizable space X is metrizable (reduce to the case where X is σ-compact by using Chapter I, § 9, no. 10, Theorem 5; then use the Corollary to Proposition 16) {[}cf.~§ 5, Exercise 15{]}.
\end{enumerate}

¶ 24) a) Let X be a complete metric space, and let A be a subset of X which is the intersection of a countable family of open sets. Show that there exists a metric on A with respect to which A is a complete metric space, and which defines on A the topology induced by that of X and also defines a uniformity finer than that induced by the uniformity of X {[}cf.~§ 1, Exercise 16 b) and c){]}.

\begin{enumerate}
\def\labelenumi{\alph{enumi})}
\setcounter{enumi}{1}
\tightlist
\item
  Conversely, let X be a metrizable space and let A be a subset of X such that there exists a metric d on A which is compatible with the topology induced by that of X and for which A is a complete metric space. Show that A is a countable intersection of open sets in X. {[}For each integer n \textgreater{} 0, consider the set \(G_{n}\) of points \(x \in \overline{A}\) which have an open neighbourhood U such that the diameter of \(A \cap U\) (with respect to d) is \(\leqslant r/n\) .{]}
\end{enumerate}

¶ 25) Let X be a metrizable space, let βX be its Stone-Čech compactification (§ 1, Exercise 7) and let d be a bounded metric compatible with the topology of X. For each x ∈ X, let f\_x(y) denote the real-valued function obtained by extending y → d(x, y) by continuity to βX.

\begin{enumerate}
\def\labelenumi{\alph{enumi})}
\item
  Show that \(f_{x}(z) + f_{y}(z) \geqslant d(x, y)\) whenever \(x\) and \(y\) are in \(\mathbf{X}\) and \(z\) is in \(\beta \mathbf{X}\) . For each \(y \in \beta \mathbf{X}\) other than \(x \in \beta \mathbf{X}\) , show that \(f_{x}(y) > 0\) . b) Show that if \(\mathbf{X}\) is complete with respect to the metric \(d\) , then \(\mathbf{X}\) is the intersection of a sequence of open sets in \(\beta \mathbf{X}\) . {[}Consider the set \(G_{n}\) of all \(y \in \beta \mathbf{X}\) such that \(f_{x}(y) < 1 / n\) for at least one point \(x \in \mathbf{X}\) , and show by using a) that \(\mathbf{X} = \bigcap_{n=1}^{\infty} G_{n}\) .{]}
\item
  Conversely, show that if X is the intersection of a sequence of open sets \(G_{n}\) in \(\beta X\) , there exists a metric \(d'\) on X which is compatible with the topology of X, and with respect to which X is complete. {[}Considering only the case in which \(X \neq \beta X\) , note that \(\beta X - X\) is the union of a family of compact sets \(F_{n} = \beta X - G_{n}\) ; for each n, let \(g_{n}(x) = \inf_{y \in F_{n}} f_{x}(y)\) , and show by using a) that \(g_{n}(x) > 0\) for all \(x \in X\) , and that \(g_{n}\) is continuous on X. Finish the proof as in Exercise 16 b) of § 1.{]}
\end{enumerate}

\ExerciseSection

\begin{enumerate}
\def\labelenumi{\arabic{enumi})}
\tightlist
\item
  A pseudometric \(f\) on a group \(G\) (written multiplicatively) is said to be left-invariant (resp. right-invariant) if \(f(zx, zy) = f(x, y)\) {[}resp. \(f(xz, yz) = f(x, y)\) {]} for all \(x, y, z\) in \(G\) .
\end{enumerate}

\begin{enumerate}
\def\labelenumi{\alph{enumi})}
\tightlist
\item
  If \(f\) is a left-invariant pseudometric, the real-valued function \(g(x) = f(e, x)\) on \(\mathbf{G}\) ( \(e\) being the identity element of \(\mathbf{G}\) ) satisfies the following conditions:
\end{enumerate}

\begin{enumerate}
\def\labelenumi{(\roman{enumi})}
\item
  \(g(x)\geqslant \mathrm{0}\) for all \(x\in \mathbf{G}\) , and \(g(e) = \mathrm{0};\)
\item
  \(g(x^{-1}) = g(x)\) ;
\item
  \(g(xy) \leqslant g(x) + g(y)\) .
\end{enumerate}

Conversely, if \(g\) is any real-valued function on \(G\) satisfying these conditions, then \(f(x, y) = g(x^{-1}y)\) is a left-invariant pseudometric on \(G\) . \(b)\) The topology \(\mathfrak{G}\) defined by a saturated family \((f_t)\) of left-invariant pseudometrics on a group \(G\) is compatible with the group structure of \(G\) if and only if, for each \(a \in G\) , each index \(t\) and each real number \(\alpha > 0\) , there exists an index \(x\) and a real number \(\beta > 0\) such that the relation \(f_x(e, x) \leqslant \beta\) implies \(f_t(e, axa^{-1}) \leqslant \alpha\) . If this condition is satisfied, the uniformity on \(G\) defined by the family \(f_t\) is the same as the left uniformity of the topological group obtained by endowing \(G\) with the topology \(\mathfrak{G}\) .

\begin{enumerate}
\def\labelenumi{\alph{enumi})}
\setcounter{enumi}{2}
\tightlist
\item
  On every topological group G there exists a family of left-invariant pseudometrics such that the uniformity defined by this family is the same as the left uniformity of G.
\end{enumerate}

\begin{enumerate}
\def\labelenumi{\arabic{enumi})}
\setcounter{enumi}{1}
\item
  Let \(\mathbf{G}\) be a topological group whose left and right uniformities coincide. Show that this unique uniformity can be defined by a family of pseudometrics which are simultaneously left-and right-invariant (using Exercise 3 of Chapter III, § 3, show that if \(\mathbf{V}\) is any neighbourhood of the identity element \(e\) of \(\mathbf{G}\) , then \(\mathbf{V}_0 = \bigcap_{x\in \mathbf{G}}x\mathbf{V}x^{-1}\) is a neighbourhood of \(e\) ).
\item
  Let G be a topological group, let \((f_{t})\) be a saturated family of left-invariant pseudometrics on G which define the left uniformity of G, and let \(g_{t}(x)=f_{t}(e,x)\) . Let H be a normal closed subgroup of G, and for each coset \(\dot{x}\in G/H\) let \(h_{t}(\dot{x})=\inf_{x\in \dot{x}}g_{t}(x)\) ; show that, if \(\bar{f}_{t}(\dot{x},\dot{y})=h_{t}(\dot{x}^{-1}\dot{y})\) , the \(\bar{f}_{t}\) form a family of left-invariant pseudometrics on G/H which define the left uniformity of G/H (argue as in Remark 2 of no. 1).
\item
  Show that the topology of a valued division ring is locally retrobounded (Chapter III, § 6, Exercise 22).
\end{enumerate}

¶ 5) Let \(\omega\) be an absolute value on the field \(\mathbf{Q}\) of rational numbers. a) Show that \(\omega\) is uniquely determined by its values at the prime numbers. b) If there is a prime number \(p\) such that \(\omega(p) \leqslant 1\) , show that \(\omega(q) \leqslant 1\) for every other prime number \(q\) {[}find an upper bound for \(\omega(q^n)\) by writing \(q^n\) in the scale of \(p\) , and let \(n \to \infty\) {]}.

\begin{enumerate}
\def\labelenumi{\alph{enumi})}
\setcounter{enumi}{2}
\tightlist
\item
  If there exists a prime number p such that \(\omega(p) < 1\) , show that \(\omega(q) = 1\) for every other prime number q {[}show that we cannot have \(\omega(q) < 1\) , by using the fact that for each integer n \textgreater{} 0 there exist two rational integers r and s such that \(1 = rp^{n} + sq^{n}\) {]} . Hence show that \(\omega\) is then an absolute value equivalent to the p-adic absolute value on Q. d) If \(\omega(p) > 1\) for every prime number p, show that if p and q are any two prime numbers, then
\end{enumerate}

\[
\frac {\log \omega (p)}{\log p} = \frac {\log \omega (q)}{\log q}
\]

{[}same method as in \(b\) {]}. Hence show that in this case \(\omega(x) = |x|^{\rho}\) for some \(\rho \leqslant 1\) .

\begin{enumerate}
\def\labelenumi{\arabic{enumi})}
\setcounter{enumi}{5}
\item
  Let E be a left vector space over a division ring K, with a countable base \((a_{n})\) . Let \((r_{n})\) be a decreasing sequence of numbers \textgreater0, tending to 0. For each \(x = \sum_{k} t_{k} a_{k} \neq 0\) of E, put \(\|x\| = r_{h}\) where h is the smallest of the indices k such that \(t_{k} \neq 0\) , and put \(\|0\| = 0\) . Show that \(\|x - y\|\) is an invariant metric on the additive group E, and that the topology it defines on E is independent of the sequence \((r_{n})\) (decreasing and tending to 0) chosen. Hence show that if the definitions of a norm and a normed space (Definitions 5 and 6) are extended to the case where the absolute value on the scalars is improper, Proposition 7 and Theorem 1 are no longer valid.
\item
  Let E be a normed space over a valued division ring. Show that if every absolutely summable family in E is summable in E, then E is complete. {[}Let \((x_{n})\) be a Cauchy sequence in E, and consider a subsequence \((x_{n_{k}})\) of \((x_{n})\) such that the series whose general term is \(x_{n_{k+1}} - x_{n}\) is absolutely convergent.{]}
\item
  Give an example of a family which is summable but not absolutely summable in the field \(Q_{p}\) of p-adic numbers. {[}Show that a family \((x_{t})_{t\in\mathbf{I}}\) of points of \(Q_{p}\) is summable if and only if \(\lim x_{t}=0\) with respect to the filter of complements of finite subsets of I.{]}
\item
  A mapping \(w\) of a ring \(\mathbf{A}\) into \(\mathbf{R}_+\) is called a semi-absolute value on \(\mathbf{A}\) if it satisfies the conditions: (i) \(w(\mathrm{0}) = \mathrm{0}\) ; (ii) \(w(x - y) \leqslant w(x) + w(y)\) ;
\end{enumerate}

\begin{enumerate}
\def\labelenumi{(\roman{enumi})}
\setcounter{enumi}{2}
\tightlist
\item
  \(w(xy) \leqslant w(x)w(y)\) . The semi-absolute value w is said to be Hausdorff if \(w(x) = 0\) implies x = 0. If w is a Hausdorff semi-absolute value on A, then \(w(x - y)\) is an invariant metric on the additive group of A, and hence defines a topology compatible with the additive group structure of A. Generalize the principal properties of normed algebras, notably Proposition 13, to rings endowed with a semi-absolute value.
\end{enumerate}

If \(w\) is a non-Hausdorff semi-absolute value on \(A\) , the set \(\overline{w}(0)\) is a two-sided ideal \(\mathfrak{a}\) in \(A\) ; if we endow \(A\) with the topology defined by the pseudometric \(w(x - y)\) , this topology is compatible with the ring structure of \(A\) , and the Hausdorff space associated with \(A\) is the quotient ring \(A / \mathfrak{a}\) , on which the function \(\overline{w}(\overline{x})\) , which for each \(\overline{x} \mod \mathfrak{a}\) is equal to the common value of \(w(x)\) for all \(x \in \overline{x}\) , is a Hausdorff semi-absolute value said to be associated with \(w\) ; the topology defined by \(\overline{w}\) is the quotient by \(\mathfrak{a}\) of the topology of \(A\) .

¶ 10) a) Two semi-absolute values \(w_1, w_2\) on a ring A are said to be equivalent if the pseudometrics \(w_1(x - y), w_2(x - y)\) are equivalent. Show that, if \(w\) is a semi-absolute value on A, then \(aw\) and \(w^{1/a}\) are semi-absolute values equivalent to \(w\) for every real number \(a \geqslant 1\) .

\begin{enumerate}
\def\labelenumi{\alph{enumi})}
\setcounter{enumi}{1}
\tightlist
\item
  If \(w_{i}\) ( \(1 \leqslant i \leqslant n\) ) are semi-absolute values on a ring A, the functions \(w = \sum_{i} w_{i}\) and \(w' = \sup_{i} w_{i}\) are two equivalent semi-absolute values on A. If
\end{enumerate}

\[
\mathfrak {a} _ {i} = \overline {{{w}}} _ {i} ^ {1} (\mathrm{0}) \quad \text { and } \quad \mathfrak {a} = \overline {{{w}}} ^ {1} (\mathrm{0}),
\]

then \(\mathfrak{a} = \bigcap_{i} \mathfrak{a}_{i}\) . If \(A_i\) denotes the completion of the quotient ring \(A/a_i\) endowed with the Hausdorff semi-absolute value associated with \(w_i\) , show that \(A/a_i\) endowed with the Hausdorff semi-absolute value associated with \(w\) is isomorphic to a subring of the product ring \(\prod_{i} A_i\) ; show that (assuming that \(A\) has an identity element \(1\) ) this ring is dense in \(\prod_{i} A_i\) if and only if the \(w_i\) satisfy the following condition: for each \(\varepsilon > 0\) and for each index \(i\) , there is an element \(x_i \in A\) such that \(w_i(1 - x_i) \leqslant \varepsilon\) and \(w_k(x_i) \leqslant \varepsilon\) for each index \(k \neq i\) .

\begin{enumerate}
\def\labelenumi{\arabic{enumi})}
\setcounter{enumi}{10}
\tightlist
\item
  Let A be a ring endowed with a Hausdorff semi-absolute value, and suppose that A has an identity element and is complete with respect to the topology defined by the semi-absolute value. Show that every maximal ideal of A is closed (use Proposition 13).
\end{enumerate}

¶ 12) Let K be a non-discrete, Hausdorff, topological division ring, and let \(\varphi: \mathbf{K} \to \mathbf{R}_{+}\) be a mapping such that \(\varphi(0) = 0\) , \(\varphi(xy) = \varphi(x)\varphi(y)\) and such that if \(V_{n}\) denotes the set of all \(x \in K\) such that \(\varphi(x) \leqslant 1/n\) , the sets \(V_{n}\) form a fundamental system of neighbourhoods of \(0\) in K.

\begin{enumerate}
\def\labelenumi{\alph{enumi})}
\tightlist
\item
  Show that there exists a real number \(a > 0\) such that, for all \(x \in \mathbf{K}\) ,
\end{enumerate}

\[
\varphi (1 + x) \leqslant a (1 + \varphi (x))
\]

{[}if not, there would exist a sequence \((x_{n})\) of points of K such that both \((1 + x_n)^{-1}\) and \(x_{n}(1 + x_{n})^{-1}\) tended to zero{]}. Hence show that if \(\psi (x) = (\varphi (x))^{\alpha}\) we have

\[
\psi (x + y) \leqslant 2 \sup (\psi (x), \psi (y))
\]

for \(\alpha\) sufficiently small.

b) Show that if \(n = 2^p\) we have \(\psi\left(\sum_{i=1}^{n} x_i\right) \leqslant n \sup (\psi(x_i))\) , and deduce that for every integer \(m > 0\) we have \(\psi(m) \leqslant 2m\) .

\begin{enumerate}
\def\labelenumi{\alph{enumi})}
\setcounter{enumi}{2}
\tightlist
\item
  Deduce from b) that for every \(n = 2^p\) and every \(x \in \mathbf{K}\) we have \(\psi((1 + x)^{n-1}) \leqslant 2n(1 + \psi(x))^{n-1}\) , and hence that \(\psi\) is an absolute value on \(\mathbf{K}\) which defines the topology of \(\mathbf{K}\) .
\end{enumerate}

\begin{enumerate}
\def\labelenumi{\arabic{enumi})}
\setcounter{enumi}{12}
\tightlist
\item
  Let K be a non-discrete Hausdorff topological division ring. Let R be the set of all \(x \in K\) such that \(\lim_{n \to \infty} x^n = 0\) , and let N be the complement of the set \(R \cup R^{-1}\) . Show that there exists an absolute value on K which defines the topology of K if and only if (i) R is open in K; (ii) for each neighbourhood V of 0 in K, there is a neighbourhood U of 0 in K such that \(R.U \subset V\) ; and (iii) if \(x \in R\) and \(y \in R \cup N\) , then \(yx \in R\) .
\end{enumerate}

To show that these conditions are sufficient, prove successively that: a) N is a normal subgroup of the multiplicative group \(K^{*}\) of non-zero elements of K.

\begin{enumerate}
\def\labelenumi{\alph{enumi})}
\setcounter{enumi}{1}
\item
  In the quotient group \(K^{*}/N\) , put \(\dot{x} \leqslant \dot{y}\) if there exist \(x \in \dot{x}\) and \(y \in \dot{y}\) such that \(yx^{-1} \in R \cup N\) ; show that this relation is an ordering compatible with the group structure of \(K^{*}/N\) , and that the ordered group so defined is isomorphic to a subgroup of the additive group R (use Exercise 1 of Chapter V, § 3).
\item
  Complete the proof with the help of Exercise 12.
\end{enumerate}

In particular, the topology of a non-discrete locally compact topological field can be defined by an absolute value.

\ExerciseSection

\begin{enumerate}
\def\labelenumi{\roman{enumi})}
\tightlist
\item
  \begin{enumerate}
  \def\labelenumii{\alph{enumii})}
  \tightlist
  \item
    Construct a topological space consisting of four points which satisfies axiom \((\mathbf{O}_{\mathbf{V}})\) but not axiom \((\mathbf{O}_{\mathbf{III}})\) .
  \end{enumerate}
\end{enumerate}

\begin{enumerate}
\def\labelenumi{\alph{enumi})}
\setcounter{enumi}{1}
\item
  A topological space X which satisfies axiom \((\mathrm{O}_{\mathrm{V}})\) also satisfies axiom \((\mathrm{O}_{\mathrm{III}})\) if and only if it has the following property: every closed subset of X is the intersection of its neighbourhoods. X is then uniformizable and the completely regular space associated with X (§ 1, Exercise 4) is normal.
\item
  If a topological space satisfies axioms (C) and \((\mathrm{O}_{\mathrm{III}})\) , it satisfies \((\mathrm{O}_{\mathrm{V}})\) , and the associated completely regular space is compact.
\end{enumerate}

\begin{enumerate}
\def\labelenumi{\arabic{enumi})}
\setcounter{enumi}{1}
\tightlist
\item
  Let X be a topological space which satisfies axiom \((\mathbf{O}_{\mathbf{v}})\) .
\end{enumerate}

\begin{enumerate}
\def\labelenumi{\alph{enumi})}
\item
  Show that the relation \(R: \{x\} \cap \{y\} \neq \emptyset\) between two points x, y of X is equivalent to the relation \(R' : "f(x) = f(y)\) for every real-valued continuous function f on X''; hence show that R is an equivalence relation on X.
\item
  Show that the quotient space \(\mathbf{Y} = \mathbf{X} / \mathbf{R}\) is normal and that, if \(\varphi: \mathbf{X} \to \mathbf{Y}\) is the canonical mapping, every continuous real-valued function \(f\) on \(\mathbf{X}\) can be written in the form \(f = g \circ \varphi\) , where \(g\) is a continuous real-valued function on \(\mathbf{Y}\) .
\end{enumerate}

\begin{enumerate}
\def\labelenumi{\arabic{enumi})}
\setcounter{enumi}{2}
\tightlist
\item
  If X is a topological space, the following conditions are equivalent:
\end{enumerate}

\begin{enumerate}
\def\labelenumi{\alph{enumi})}
\item
  Every subspace of X satisfies axiom \((\mathrm{O}_{\mathrm{V}}^{\prime})\) .
\item
  Every open subspace of X satisfies axiom \((\mathbf{O}_{\mathbf{V}}^{\prime})\) .
\item
  Given any two subsets A and B of X such that \(A \cap \overline{B} = B \cap \overline{A} = \varnothing\) , there exist two disjoint open sets U, V such that \(A \subset U\) and \(B \subset V\) .
\end{enumerate}

A topological space \(\mathbf{X}^{-}\) is said to be completely normal if it is Hausdorff and satisfies these conditions. Every metrizable space is completely normal. A compact space need not be completely normal.

\begin{enumerate}
\def\labelenumi{\arabic{enumi})}
\setcounter{enumi}{3}
\tightlist
\item
  Every linearly ordered set X endowed with either of the topologies \(\mathcal{C}_{+}(\mathbf{X})\) , \(\mathcal{C}_{-}(\mathbf{X})\) (Chapter I, § 2, Exercise 5) is completely normal. (If A and B are two subsets of X such that \(\mathbf{A} \cap \overline{\mathbf{B}} = \mathbf{B} \cap \overline{\mathbf{A}} = \emptyset\) , define a neighbourhood \(V_{x}\) of \(x\) for each \(x \in \mathbf{A}\) , and a neighbourhood \(W_{y}\) of \(y\) for each \(y \in \mathbf{B}\) , such that \(V_{x} \cap W_{y} = \emptyset\) for all \(x \in \mathbf{A}\) and all \(y \in \mathbf{B}\) .)
\end{enumerate}

¶ 5) Every linearly ordered set X, endowed with the topology \(\mathcal{G}_{0}(\mathbf{X})\) (Chapter I, § 2, Exercise 5), is completely normal. {[}Consider first the case where X is compact; show with the help of Chapter IV, § 2, Exercise 6 that every open set in X is the union of mutually disjoint open intervals; use this result to prove the proposition, by considering (in the notation of Exercise 3) the complement of the closed set \(\overline{\mathbf{A}}\cap\overline{\mathbf{B}}\) , and then the complement of \(\overline{\mathbf{B}}\) ; to pass to the general case, in which X is arbitrary, use Chapter IV, § 4, Exercise 7.{]}

\begin{enumerate}
\def\labelenumi{\arabic{enumi})}
\setcounter{enumi}{5}
\item
  Show that, in a normal space X, every subspace Y of X which is a countable union of closed sets is normal. {[}Let A, B be two disjoint closed subsets of Y, and suppose that Y is the union of an ascending sequence \((\mathbf{Y}_n)\) of closed subsets of \(\mathbf{X}\) . Define by induction two sequences \((\mathbf{U}_n), (\mathbf{V}_n)\) of open subsets of \(\mathbf{X}\) such that: (i) \(\overline{\mathbf{U}_n \cap \mathbf{Y}_n}\) and \(\overline{\mathbf{V}_n \cap \mathbf{Y}_n}\) do not intersect; (ii) \(\mathbf{U}_n \cap \mathbf{Y}_n\) contains \(\mathbf{A} \cap \mathbf{Y}_n\) and all the sets \(\overline{\mathbf{U}_i \cap \mathbf{Y}_i}\) for \(1 \leqslant i \leqslant n\) ; and \(\mathbf{V}_n \cap \mathbf{Y}_n\) contains \(\mathbf{B} \cap \mathbf{Y}_n\) and all \(\overline{\mathbf{V}_i \cap \mathbf{Y}_i}, 1 \leqslant i \leqslant n\) .{]}
\item
  A topological space X is said to be perfectly normal if it is normal and if each closed subset of X is a countable intersection of open sets (or, equivalently, if each open subset of X is a countable union of closed sets).
\end{enumerate}

\begin{enumerate}
\def\labelenumi{\alph{enumi})}
\item
  A Hausdorff space X is perfectly normal if and only if, given any closed subset A of X, there exists a real-valued continuous function f on X such that \(\overline{f}^{1}(0)=A\) {[}use the method of Exercise 16 a) of §1{]}. b) Show that every perfectly normal space X is completely normal and that every subspace of X is perfectly normal {[}use Exercises 3 b) and 6{]}.
\item
  Every lower semi-continuous real-valued function on a perfectly normal space is the upper envelope of a sequence of continuous functions (use the method of § 2, no. 7, Proposition 11).
\item
  The compact space X obtained by adjoining a point at infinity to an uncountable discrete space is completely normal but not perfectly normal, and every subspace of X is paracompact.
\end{enumerate}

§ 8) Show that the non-metrizable compact space X defined in Exercise 13 a) of § 2 is perfectly normal, and that the non-metrizable compact space Y defined in Exercise 13 d) of § 2 is completely normal but not perfectly normal.

¶ 9) a) Let X be a paracompact space in which every point has a countable fundamental system of neighbourhoods, and let Y be a countably compact normal space (§ 2, Exercise 14). Show that X × Y is normal. {[}Let A, B be two disjoint closed subsets of X × Y; for each x ∈ X, show that there is a neighbourhood U\_x of x in X and two open sets V\_x, W\_x in Y such that \(\overline{V}_x \cap \overline{W}_x = \emptyset\) and such that for each z ∈ U\_x we have A(z) ⊂ V\_x and B(z) ⊂ W\_x; to prove this, argue by contradiction. Let (T\_λ)\_\{λ ∈ L\} be a locally finite open covering of X, finer than the covering (U\_x)\_\{x ∈ X\}, and let (f\_λ) be a partition of unity subordinate to (T\_λ); for each λ ∈ L, let x\_λ be such that T\_λ ⊂ U\_\{x\_λ\} and let g\_λ be a continuous mapping of Y into {[}0, 1{]} which is equal to 0 on \(\overline{V}_{x_λ}\) and to 1 on \(\overline{W}_{x_λ}\) ; consider the function \(\sum_{\lambda} f_\lambda(x) g_\lambda(y)\) on X × Y.{]}

\begin{enumerate}
\def\labelenumi{\alph{enumi})}
\setcounter{enumi}{1}
\tightlist
\item
  Let Y = {[}a, b{[} be the locally compact space defined in Chapter I, § 9, Exercise 12, and let X = {[}a, b{]} be the compact space obtained by adjoining a point at infinity to Y. Y is normal (Exercise 4) and countably compact (§ 2, Exercise 15), but the product X×Y is not normal (use Chapter II, § 4, Exercise 4).
\end{enumerate}

¶ 10) Let L be an uncountable set and let X denote the completely regular space \(\mathbf{N}^{\mathrm{L}}\) (N carrying the discrete topology). Let A (resp. B) be the set of all \(x = (x(\lambda))_{\lambda \in \mathrm{L}}\) in X such that, for each integer \(k \neq 0\) (resp. \(k \neq 1\) ) the set of \(\lambda \in \mathbf{L}\) such that \(x(\lambda) = k\) has at most one element. a) Show that A and B are disjoint closed subsets of X.

\begin{enumerate}
\def\labelenumi{\alph{enumi})}
\setcounter{enumi}{1}
\item
  Let U, V be two open subsets of X such that \(A \subset U\) and \(B \subset V\) . Let \(\Phi\) denote the set of all elementary sets in X (Chapter I, § 4, no. 1) whose projections which are distinct from N consist of a single element; for each \(W \in \Phi\) , let \(\mathrm{H}(\mathrm{W})\) be the set of all \(\lambda \in L\) such that \(\mathrm{pr}_{\lambda}(\mathrm{W})\) consists of a single element. Show that there is a sequence \((x_{n})\) of points of A, a sequence \((\mathbf{U}_{n})\) of sets of \(\Phi\) , a sequence \((\lambda_{n})\) of distinct elements of L and a strictly increasing sequence \((m(n))_{n \in \mathbb{N}}\) of integers with the following properties: (i) \(U_{n} \subset U\) is a neighbourhood of \(x_{n}\) ; (ii) \(\mathrm{H}(\mathrm{U}_{n})\) is the set of all \(\lambda_{k}\) such that \(k \leqslant m(n)\) ; (iii) \(x_{0}(\lambda) = 0\) for each \(\lambda \in L\) , and for each n \textgreater{} 0, \(x_{n}(\lambda_{k}) = k\) if \(k \leqslant m(n - 1)\) and \(x_{n}(\lambda) = 0\) for all other \(\lambda \in L\) .
\item
  Let \(y \in \mathbf{B}\) be the point such that \(y(\lambda_k) = k\) for each integer \(k\) and \(y(\lambda) = \mathbf{1}\) for all \(\lambda \in \mathbf{L}\) other than the \(\lambda_k\) . Let \(\mathbf{V}_0 \subset \mathbf{V}\) be a set of \(\Phi\) which contains \(y\) , and let \(n\) be an integer such that \(\lambda_k \in \mathbf{L} - \mathbf{H}(\mathbf{V}_0)\) for all \(k > m(n)\) . Show that \(\mathbf{U}_{n+1} \cap \mathbf{V}_0 \neq \emptyset\) , and deduce that \(\mathbf{X}\) is not normal.
\end{enumerate}

\begin{enumerate}
\def\labelenumi{\arabic{enumi})}
\setcounter{enumi}{10}
\tightlist
\item
  Deduce from Exercise 10 that:
\end{enumerate}

\begin{enumerate}
\def\labelenumi{\alph{enumi})}
\item
  If a product \(\prod_{i\in I}X_i\) of non-empty Hausdorff spaces is normal, then \(X_i\) is countably compact (§ 2, Exercise 14) for all save a countable set of indices.
\item
  A product \(\prod_{i\in I}X_i\) of non-empty metrizable spaces is normal if and only if \(X_{i}\) is compact for all but a countable set of indices; and the product space \(\prod_{i\in I}X_{i}\) is then paracompact.
\end{enumerate}

\begin{enumerate}
\def\labelenumi{\arabic{enumi})}
\setcounter{enumi}{11}
\tightlist
\item
  Let X, Y be two Hausdorff spaces.
\end{enumerate}

\begin{enumerate}
\def\labelenumi{\alph{enumi})}
\item
  Suppose that there is a closed set A in X which is not a countable intersection of open sets, and that there is a countably infinite subset B of Y which is not closed. Let \(b \in \overline{B} - B\) , and consider the subsets \(C = A \times B\) , \(D = (X - A) \times \{b\}\) in \(X \times Y\) . Show that \(C \cap \overline{D} = \overline{C} \cap D = \varnothing\) , but that there is no pair U, V of open subsets of \(X \times Y\) such that \(C \subset U\) , \(D \subset V\) and \(U \cap V = \varnothing\) .
\item
  If \(X \times Y\) is completely normal then either (i) one of the spaces X, Y is perfectly normal, or (ii) one of the spaces X, Y has the property that every countably infinite subset is closed (cf.~§ 5, Exercise 16).
\item
  Let \(X = [a, b]\) be an uncountable well-ordered set such that \([a, x]\) is countable for each x \textless{} b. Endow X with the topology in which \(\{x\}\) is open for each x \textless{} b and the intervals \([x, b]\) (x \textless{} b) form a fundamental system of neighbourhoods of b. Show that \(X \times X\) is completely normal but that X is not perfectly normal.
\item
  Show that if \(\mathbf{X}\) is a compact space such that \(\mathbf{X} \times \mathbf{X} \times \mathbf{X}\) is completely normal, then \(\mathbf{X}\) is metrizable (use Theorem 1 of §2, no. 4).
\end{enumerate}

¶ 13) Let \((\mathbf{X}_{t})_{t\in I}\) be an uncountable family of Hausdorff topological spaces, each containing more than one point. For each \(t\in I\) , let \(a_{t}\) and \(b_{t}\) be two distinct points of \(\mathbf{X}_{t}\) . Let \(\mathbf{X}\) be the product space \(\prod_{t\in I}\mathbf{X}_{t}\) , and let Y be the subspace of X consisting of these points \((x_{t})_{t\in I}\) such that \(x_{t}=a_{t}\) for all but a countable set of indices t; and let b be the point \(b_{t}\) of X.

In the product space \(\mathbf{X} \times \mathbf{Y}\) , let \(\mathbf{A}\) be the diagonal of \(\mathbf{Y} \times \mathbf{Y}\) and let \(\mathbf{B}\) be the set \(\{b\} \times \mathbf{Y}\) . Show that \(\mathbf{A}\) and \(\mathbf{B}\) are closed sets and that there is no pair of open sets \(\mathbf{U}, \mathbf{V}\) in \(\mathbf{X} \times \mathbf{Y}\) such that \(\mathbf{A} \subset \mathbf{U}, \mathbf{B} \subset \mathbf{V}\) , and \(\mathbf{U} \cap \mathbf{V} = \emptyset\) . {[}Let \(\mathbf{U}, \mathbf{V}\) be two open sets such that \(\mathbf{A} \subset \mathbf{U}\) and \(\mathbf{B} \subset \mathbf{V}\) ; show that there is an increasing sequence \((\mathbf{H}_n)\) of countable subsets of \(\mathbf{I}\) , such that if \(x_n\) denotes the point of \(\mathbf{Y}\) for which each coordinate with index \(i \in \mathbf{H}_n\) is equal to \(b_i\) , and each coordinate with index \(i \notin \mathbf{H}_n\) is equal to \(a_i\) , then the point \((x_{n+1}, x_n)\) belongs to \(\mathbf{V}\) for each integer \(n\) ; prove that the sequence \((x_{n+1}, x_n)\) tends to a point of \(\mathbf{A}\) , and deduce that \(\mathbf{U} \cap \mathbf{V} \neq \emptyset\) .{]}

Hence construct an example of a connected topological ring which is not normal.

Deduce also that the product of an uncountable family of Hausdorff topological spaces, each of which has at least two distinct points, cannot be completely normal (use the result above when each \(X_{t}\) is a discrete space of two points).

\begin{enumerate}
\def\labelenumi{\arabic{enumi})}
\setcounter{enumi}{13}
\item
  Let X be a non-normal Hausdorff space, and let A, B be two disjoint closed subsets of X such that there is no pair of disjoint open sets U, V in X satisfying the relations \(A \subset U\) and \(B \subset V\) . Let R be the equivalence relation on X whose equivalence classes are the set A, the set B, and the sets \(\{x\}\) where \(x \in \complement(A \cup B)\) . Show that the graph of R is closed in \(X \times X\) and that the relation R is closed, but that the quotient space X/R is not Hausdorff (cf.~Chapter I, § 8, no. 3, Proposition 8).
\item
  \begin{enumerate}
  \def\labelenumii{\alph{enumii})}
  \tightlist
  \item
    Let X be a normal space and let R be a closed equivalence relation on X. Show that the quotient space X/R is normal (use Proposition 10 of Chapter I, § 5, no. 4).
  \end{enumerate}
\end{enumerate}

\begin{enumerate}
\def\labelenumi{\alph{enumi})}
\setcounter{enumi}{1}
\item
  Let X be a locally compact σ-compact space, and let R be an equivalence relation on X whose graph is closed in X×X. Show that X/R is normal (cf.~Chapter I, § 10, Exercise 19).
\item
  Let X be the complement in R of the set of points of the form 1/n, where n is any integer except 0 or ±1. Consider the equivalence relation R on the metrizable space X for which the equivalence class of each \(x \in X\) , which is not an integer, consists of the points x and 1/x, and the equivalence class of each integer consists of that integer alone. Show that R is open and that the graph of R is closed in \(X \times X\) , but that X/R is not regular.
\end{enumerate}

\begin{enumerate}
\def\labelenumi{\arabic{enumi})}
\setcounter{enumi}{15}
\tightlist
\item
  For each covering \(\Re\) of a topological space \(X\) , let \(V_{\Re}\) denote the union of all the sets \(U \times U\) such that \(U \in \Re\) . A covering \(\Re\) of \(X\) is said to be even if there is a neighbourhood \(W\) of the diagonal \(\Delta\) of \(X \times X\) such that the covering \((W(x))_{x \in X}\) is finer than \(\Re\) . \(\Re\) is said to be divisible if there is a neighbourhood \(W\) of \(\Delta\) in \(X \times X\) such that \(\dot{W} \subset V_{\Re}\) . a) Show that every even open covering of \(X\) is divisible {[}cf.~Exercise 19 b){]}.
\end{enumerate}

\begin{enumerate}
\def\labelenumi{\alph{enumi})}
\setcounter{enumi}{1}
\tightlist
\item
  Let \(\Re\) be an open covering of \(X\) such that there exists a locally finite closed covering \(\mathfrak{S}\) of \(X\) which is finer than \(\Re\) . Show that \(\Re\) is even. (For every set \(A \in \mathfrak{S}\) , let \(U_A\) be a set of \(\Re\) which contains \(A\) , and let \(V_A\) be the union in \(X \times X\) of \(U_A \times U_A\) and \((X - A) \times (X - A)\) ; consider the set \(W = \bigcap_{A \in \mathfrak{S}} V_A\) .)
\end{enumerate}

¶ 17) a) If every finite open covering of a topological space X is divisible (Exercise 16), then X satisfies Axiom \((\mathrm{O}_{\mathbf{V}}^{\prime})\) (consider a binary open covering of X, i.e., one formed by two subsets of X). Conversely, if X satisfies \((\mathrm{O}_{\mathbf{V}}^{\prime})\) , then every finite open covering of X is even {[}use Theorem 3 of no. 3, which is valid provided that X satisfies \((\mathrm{O}_{\mathbf{V}}^{\prime})\) , to reduce to the case of a binary covering, and then again to deal with this case{]}. As R runs through the set of finite open coverings of X, the sets \(V_{R}\) form a fundamental system of entourages of a uniformity on X, compatible with the topology of X; this uniformity is called the ``uniformity of finite open coverings''.

\begin{enumerate}
\def\labelenumi{\alph{enumi})}
\setcounter{enumi}{1}
\item
  Suppose that X is normal. Show that the uniformity U of finite open coverings is the same as the coarsest uniformity \(U'\) for which all continuous mappings of X into {[}0, 1{]} are uniformly continuous {[}i.e., the uniformity induced on X by that of its Stone-Čech compactification (§ 1, Exercise 7){]}. {[}To show that U is coarser than \(U'\) , use Theorem 3 of no. 3 and Axiom \((\mathrm{O_V})\) in order to construct a finite family of pseudometrics on \(\mathbf{X}\) of the form \((x, y) \to |f(x) - f(y)|\) , in such a way that an entourage of the uniformity defined by these pseudometrics is contained in an entourage of the uniformity \(\mathfrak{U}]\) .
\item
  Let X be a normal space and let \(\beta\mathbf{X}\) be its Stone-Čech compactification. For each closed subset Y of X, let \(\overline{\mathbf{Y}}\) be the closure of Y in \(\beta\mathbf{X}\) . Show that the continuous mapping of \(\beta\mathbf{Y}\) into \(\overline{\mathbf{Y}}\) which extends the identity mapping of Y ( \(\beta\mathbf{Y}\) being the Stone-Čech compactification of Y) is a homeomorphism {[}either use \(b\) ) or else Theorem 2 and the universal property of \(\beta\mathbf{Y}]\) . If Y and Z are two closed subsets of X, show that \(\overline{\mathbf{Y}} \cap \overline{\mathbf{Z}} = \overline{\mathbf{Y} \cap \mathbf{Z}}\) in \(\beta\mathbf{X}\) (if \(x_0 \notin \overline{\mathbf{Y}} \cap \overline{\mathbf{Z}}\) and \(x_0 \in \overline{\mathbf{Y}}\) , consider a closed neighbourhood V of \(x_0\) in \(\beta\mathbf{X}\) , such that V \(\cap\) Y and Z are disjoint closed subsets of X).
\end{enumerate}

¶ 18) A family \((A_{\alpha})\) of subsets of a topological space X is said to be discrete if, for each \(x \in X\) , there is a neighbourhood of x which meets at most one set \(A_{\alpha}\) of the family. A Hausdorff space X is said to be collectively normal if, for each discrete family \((A_{\alpha})\) of closed subsets of X, there is a family \((U_{\alpha})\) of mutually disjoint open subsets of X such that \(A_{\alpha} \subset U_{\alpha}\) for each index \(\alpha\) . Every collectively normal space is normal.

\begin{enumerate}
\def\labelenumi{\alph{enumi})}
\item
  Every open covering of a completely regular space X is divisible (Exercise 16) if and only if the set of neighbourhoods of the diagonal \(\Delta\) in \(\mathbf{X} \times \mathbf{X}\) is the filter of entourages of the universal uniformity of X ( \(\S\) 1, Exercise 5). Show that if this condition is satisfied, then X is collectively normal {[}if \((\mathbf{A}_{\alpha})\) is a discrete family of closed subsets of X, consider the open covering \((\mathbf{V}_{\alpha})\) , where \(\mathbf{V}_{\alpha} = \mathbf{X} - \bigcup_{\beta \neq \alpha} \mathbf{A}_{\beta}\) for each \(\alpha\) {]}.
\item
  Let Y = {[}a, b{]} be the locally compact space defined in Chapter I, § 9, Exercise 12; let Y₀ be the set Y endowed with the discrete topology, and let Z be the set Y₀ ∪ \{b\}, where the points of Y₀ are open sets, and the sets {]}x, b{[}, where x runs through Y₀, form a fundamental system of neighbourhoods of b. Show that the space X = Y × Z is collectively normal (recall that no infinite family of subsets of Y can be discrete, and use Exercise 4). Let R be the open covering of X consisting of Y × Y₀ and the products {[}a, x{]} × {]}x, b{]} (where x \textless{} b); show that R is not divisible and that consequently the set of neighbourhoods of the diagonal Δ in X × X is not the filter of entourages of a uniformity on X {[}use Exercise 12 a) of Chapter I, § 9{]}.
\end{enumerate}

¶ 19) a) A regular space X is paracompact if and only if every open covering of X is even (Exercise 16). {[}To prove necessity, use Lemma 5 and Exercise 16 b); for sufficiency, use Exercise 16 a) of § 4, Proposition 2 of § 1, no. 4, and Theorem 4 of § 4, no. 5{]}.

\begin{enumerate}
\def\labelenumi{\alph{enumi})}
\setcounter{enumi}{1}
\item
  Give an example of a divisible non-even covering of a collectively normal, non-paracompact space {[}use a) and Chapter II, § 4, Exercise 4{]}.
\item
  Show that a paracompact space X is complete with respect to its universal uniformity (§ 1, Exercise 5). {[}If a Cauchy filter F with respect to this uniformity has no cluster point, then every point \(x \in X\) has an open neighbourhood \(V_{x}\) which does not meet at least one set of F; use a) and Exercise 18 a).{]}
\end{enumerate}

\begin{enumerate}
\def\labelenumi{\arabic{enumi})}
\setcounter{enumi}{19}
\tightlist
\item
  \begin{enumerate}
  \def\labelenumii{\alph{enumii})}
  \tightlist
  \item
    A regular space X is paracompact if, for each open covering R of X, there is a sequence \((\mathfrak{S}_{n})\) such that \(\mathfrak{S} = \bigcup_{n} \mathfrak{S}_{n}\) is an open covering of \(\mathbf{X}\) which refines \(R\) and such that each \(\mathfrak{S}_n\) is a locally finite family (cf.~Lemmas 4, 5 and 6).
  \end{enumerate}
\end{enumerate}

\begin{enumerate}
\def\labelenumi{\alph{enumi})}
\setcounter{enumi}{1}
\item
  Deduce from a) that in a paracompact space every countable union of closed sets is a paracompact subspace (cf.~§ 5, Exercise 15).
\item
  In a regular space X, let F be a locally finite family of closed sets each of which is a paracompact subspace of X. Show that the union of the sets of F is a paracompact subspace of X (use Lemma 6 of no. 5; cf.~§ 5, Exercise 15).
\item
  Show that if X is a paracompact space and Y is a regular space which is a countable union of compact sets, then \(X \times Y\) is paracompact {[}embed Y in its Stone-Čech compactification and use b); cf.~§ 5, Exercise 16{]}.
\end{enumerate}

¶ 21) a) Let X be a paracompact space and let R be a closed equivalence relation on X such that every equivalence class with respect to R is compact. Show that X/R is paracompact {[}use Exercise 15 a), Theorem 3 of no. 3, and Lemma 6 of no. 5; cf.~Exercise 15 c){]}.

\begin{enumerate}
\def\labelenumi{\alph{enumi})}
\setcounter{enumi}{1}
\item
  Let X be a Hausdorff space and let R be a closed equivalence relation on X such that every equivalence class with respect to R is compact. Show that if X/R is paracompact, then so is X (use the argument of Chapter I, § 9, no. 10, Proposition 17).
\item
  Let Y be a locally compact but not paracompact space (cf.~Chapter I, § 9, Exercise 12), and let \(Y_{0}\) be the one-point compactification of Y. Let X be the topological sum of Y and \(Y_{0}\) , and let R be the equivalence relation on X which identifies each point of Y with its canonical image in \(Y_{0}\) . The relation R is open, every equivalence class mod R contains at most two points, and X/R is homeomorphic to \(Y_{0}\) , hence is compact.
\end{enumerate}

\begin{enumerate}
\def\labelenumi{\arabic{enumi})}
\setcounter{enumi}{21}
\item
  A regular space \(\mathbf{X}\) is metrizable if and only if there exists a sequence \((\mathfrak{B}_n)\) of locally finite families of open subsets of \(\mathbf{X}\) such that \(\mathfrak{B} = \bigcup_{n} \mathfrak{B}_n\) is a base of the topology of \(\mathbf{X}\) . {[}To show that the condition is necessary, use Lemma 3 of no. 5. To show that it is sufficient, show first that X is paracompact by using Lemmas 4, 5 and 6. For each pair of integers m, n and each \(U \in \mathfrak{B}_{m}\) , let \(U'\) be the union of the sets \(V \in \mathfrak{B}_{n}\) such that \(\overline{V} \subset U\) . Show that \(\overline{U}' \subset U\) . Let \(f_{U}: X \to [0, 1]\) be a continuous mapping which is equal to 0 on X---U and to 1 on \(U'\) , and let \(d_{mn}(x, y) = \sum_{U \in \mathfrak{B}_{m}} |f_{U}(x) - f_{U}(y)|\) ; show that the pseudometrics \(d_{mn}\) define the topology of X.{]} (``Theorem of Nagata-Smirnov.''). In particular, every regular space with a countable base is metrizable.
\item
  \begin{enumerate}
  \def\labelenumii{\alph{enumii})}
  \tightlist
  \item
    Every regular Lindelöf space (Chapter I, § 9, Exercise 15) is paracompact {[}use Exercise 20 a){]}.
  \end{enumerate}
\end{enumerate}

\begin{enumerate}
\def\labelenumi{\alph{enumi})}
\setcounter{enumi}{1}
\tightlist
\item
  The product of a Lindelöf space and a compact space is a Lindelöf space (argue as in the proof of Proposition 17 of Chapter I, § 9, no. 10; cf.~§ 5, Exercise 16).
\end{enumerate}

\begin{enumerate}
\def\labelenumi{\arabic{enumi})}
\setcounter{enumi}{23}
\tightlist
\item
  \begin{enumerate}
  \def\labelenumii{\alph{enumii})}
  \tightlist
  \item
    Let X be a metrizable space and let R be a closed equivalence relation on X, all of whose equivalence classes are compact. Show that X/R is metrizable. {[}Apply the Nagata-Smirnov theorem (Exercise 22) by proving the analogue of Lemma 3 of no. 5 for open coverings of X formed of sets which are saturated with respect to R; let \(F_{n\alpha}\) denote the largest saturated open set contained in the set of all \(x \in X\) such that \(d(x, X - U_{\alpha}) > 2^{-n}\) , let \(F_{n\alpha}'\) denote the saturation of the set of all \(x \in X\) such that \(d(x, X - U_{\alpha}) \geqslant 2^{-n}\) and let \(G_{n\alpha}\) denote the set of all \(x \in F_{n\alpha}\) such that \(x \notin F_{n+1,\beta}'\) for all \(\beta < \alpha\) .{]}
  \end{enumerate}
\end{enumerate}

\begin{enumerate}
\def\labelenumi{\alph{enumi})}
\setcounter{enumi}{1}
\item
  Extend the result of a) to the case in which the relation R is closed and every point of X/R has a countable fundamental system of neighbourhoods {[}use Exercise 20 b) of § 2{]}.
\item
  Let X be a topological space and let \((\mathbf{F}_{\alpha})\) be a locally finite closed covering of X. Show that if each of the subspaces \(F_{\alpha}\) is metrizable then X is metrizable (consider X as a quotient space of the topological sum of the \(F_{\alpha}\) ).
\item
  A paracompact space, in which every point has a metrizable neighbourhood, is metrizable {[}apply c); cf.~Chapter I, § 9, Exercise 12 b), also Chapter IX, § 4, Exercise 25 d), § 2, Exercise 15 and § 5, Exercise 15{]}.
\end{enumerate}

\begin{enumerate}
\def\labelenumi{\arabic{enumi})}
\setcounter{enumi}{24}
\tightlist
\item
  A Hausdorff space X is said to be metacompact if, given any open covering R of X, there is a point-finite open covering of X which refines R.
\end{enumerate}

\begin{enumerate}
\def\labelenumi{\alph{enumi})}
\item
  Show that every closed subspace of a metacompact space is metacompact, and that if every open subspace of a metacompact space X is metacompact, then every subspace of X is metacompact.
\item
  Let X be a Hausdorff space and let R be a closed equivalence relation on X, all of whose equivalence classes are compact. Show that if X/R is metacompact, then X is metacompact {[}argue as in Exercise 21 b){]}.
\item
  Show that a topological space which is both metacompact and countably compact (§ 2, Exercise 14) is compact. {[}Use Zorn's lemma to show that every point-finite open covering contains a minimal open covering, and use § 2, Exercise 14 a).{]}
\end{enumerate}

\(d)\) The locally compact, completely normal space \(\mathbf{X} = [a, b[\) defined in Chapter I, § 9, Exercise 12 is not metacompact (*).

\begin{enumerate}
\def\labelenumi{\alph{enumi})}
\setcounter{enumi}{4}
\tightlist
\item
  Let Y be a metrizable space and let A be a subset of Y such that both A and its complement are dense in Y. Let X denote the non-regular space obtained by endowing Y with the topology generated by the open sets of Y and the set A. Show that X is metacompact. (For each \(x \in X\) , consider a neighbourhood \(V_{x}\) of x contained in a set of a given open covering R of X, such that \(V_{x} \subset A\) if \(x \in A\) , and such that \(V_{x}\) is a neighbourhood of x in Y if \(x \notin A\) . Note that the subspace A of X is paracompact, and that the union of the \(V_{x}\) for \(x \notin A\) is also paracompact.)
\end{enumerate}

\begin{enumerate}
\def\labelenumi{\arabic{enumi})}
\setcounter{enumi}{25}
\tightlist
\item
  \begin{enumerate}
  \def\labelenumii{\alph{enumii})}
  \tightlist
  \item
    Show that every pseudo-compact normal space X (§ 1, Exercise 22) is countably compact (§ 2, Exercise 14). {[}Show that if \((x_n)\) is a sequence of distinct points of X with no cluster point, then there is a continuous finite real-valued function f on X such that \(f(x_n) = n\) for all n.{]}
  \end{enumerate}
\end{enumerate}

\begin{enumerate}
\def\labelenumi{\alph{enumi})}
\setcounter{enumi}{1}
\tightlist
\item
  Let \(Y = [a, b[\) be the locally compact space defined in Chapter I, § 9, Exercise 12, and let \(Y_{0}\) be the compact space obtained by adjoining to Y a point at infinity (which may be identified with b). Let c be the smallest element of Y such that {[}a, c{[} is infinite, and let Z be the compact subspace {[}a, c{]} of Y. Let X be the complement of the point (b, c) in the compact product space \(Y_{0} \times Z\) . Show that X is locally compact but not normal {[}cf.~Exercise 12 a){]}; that X is pseudo-compact but not countably compact. Give an example of a closed subspace of X which is not pseudo-compact. Also, give an example of a lower semi-continuous function on X which does not attain its greatest lower bound {[}cf.~§ 2, Exercise 14 i){]}.
\end{enumerate}

\begin{enumerate}
\def\labelenumi{\arabic{enumi})}
\setcounter{enumi}{26}
\tightlist
\item
  A topological space X is said to be locally paracompact if every point of X has a closed neighbourhood which is a paracompact subspace of X.
\end{enumerate}

\begin{enumerate}
\def\labelenumi{\alph{enumi})}
\item
  Show that a locally paracompact space is completely regular.
\item
  In the compact space \(Y_{0} \times Z\) defined in Exercise 26 b), let H be the complement of the set of points \((b, y)\) where y \textless{} c.~Show that H is normal but not locally paracompact {[}cf.~Chapter I, § 9, Exercise 12 b){]}.
\item
  Let X be a topological space and let \(\Phi\) be a covering of X by closed sets which are paracompact subspaces of X, such that each \(A \in \Phi\) has a neighbourhood in X which belongs to \(\Phi\) . Let \(X'\) be the set which is the sum of X and a point \(\omega\) ; define a topology on \(X'\) by taking as a fundamental system of neighbourhoods of \(x \in X\) in \(X'\) the set of all neighbourhoods of x in X, and as a fundamental system of neighbourhoods of \(\omega\) the filter base generated by the complements in \(X'\) of the sets of \(\Phi\) {[}cf.~Exercise 20 c){]}. Show that \(X'\) is paracompact.
\end{enumerate}

\begin{enumerate}
\def\labelenumi{\arabic{enumi})}
\setcounter{enumi}{27}
\tightlist
\item
  Let X be a metric space, let d be the metric on X, let A be a closed subset of X and let \(f: A \to [1, 2]\) be a continuous real-valued function. Show that the real-valued function g on X which is such that \(g(x) = f(x)\) if \(x \in A\) , and
\end{enumerate}

\[
g (x) = \frac {\mathrm{1}}{d (x , \mathrm{A})} \inf _ {y \in \mathrm{A}} (f (y) d (x, y))
\]

if \(x \in \mathbb{C}A\) , is continuous on \(X\) . Hence give a proof of Theorem 2 for metric spaces.

\begin{enumerate}
\def\labelenumi{\arabic{enumi})}
\setcounter{enumi}{28}
\item
  Let X be a Hausdorff topological space. Then X is normal if and only if, for each closed subset A of X and each real-valued function f defined on A and continuous with respect to the induced topology, there exists a continuous pseudometric d on X with respect to which f is uniformly continuous. (To show that the condition is sufficient, show that every real-valued function f defined on a closed subset A, and continuous with respect to the induced topology, can be extended to a continuous function on X; for this, consider the Hausdorff space \(\overline{X}\) associated with the uniformity defined on X by the pseudometric d.)
\item
  Let X be a normal space and let g (resp. f) be an upper (resp. lower) semi-continuous function on X, such that \(g \leqslant f\) . Show that there is a continuous real-valued function h on X such that \(g \leqslant h \leqslant f\) . {[}Reduce to the case where f and g take their values in {[}0, 1{]}. For each dyadic number \(r \in [0, 1]\) , let \(\mathrm{F}(r)\) {[}resp. \(\mathrm{G}(r)\) {]} be the set of all \(x \in X\) such that \(f(x) \leqslant r\) (resp. \(g(x) < r\) ). Define by induction a family of open sets \(\mathrm{U}(r)\) (where r runs through the set of dyadic numbers contained in {[}0, 1{]}) such that when \(r < r'\) we have \(\mathrm{F}(r) \subset \mathrm{U}(r')\) , \(\overline{\mathrm{U}(r)} \subset \mathrm{U}(r')\) and \(\overline{\mathrm{U}(r)} \subset \mathrm{G}(r')\) , by imitating the proof of Theorem 1 of no. 1; then complete the proof as in Theorem 1.{]}
\end{enumerate}
